\documentclass[output=paper,colorlinks,citecolor=brown]{langscibook}
\ChapterDOI{10.5281/zenodo.20611891}
\author{A. Krasovitsky\orcid{}\affiliation{University of Oxford \& Surrey Morphology Group, University of Surrey} and M. Kyuseva\orcid{}\affiliation{Surrey Morphology Group, University of Surrey} and M. Baerman\orcid{}\affiliation{Surrey Morphology Group, University of Surrey} and G. G. Corbett\orcid{}\affiliation{Surrey Morphology Group \& Department of Human Behavior, Ecology and Culture, Max Planck Institute for Evolutionary Anthropology, Leipzig}}


\title{Diachronic paths to uninflectedness in South Slavic}

\abstract{We argue that the ongoing loss of case in Serbian and Bulgarian dialects has led to a split system with competing co-grammars: one where nouns inflect and another where they do not. This division between \textit{forerunners} (without inflection) and \textit{underachievers} (with inflection) results in partial rules of different kinds. There are three types of splits depending on the dialect. First, a split between morphological classes\is{morphological!class}, where some retain case distinctions while others do not. Second, classes that retain inflection may still exhibit uninflectedness in specific constructions, such as partitive constructions. Third, an otherwise inflecting morphological class\is{morphological!class} may split on a semantic basis, such as animate vs. inanimate. These asymmetries in morphological change lead to the rise of partial synchronic rules, whether morphological, syntactic, or semantic. These rules persist over time and evolve as historical processes move into new phases. By studying these rules in contemporary dialects and arranging them chronologically, we can uncover the detailed stages of this historical process.}


\IfFileExists{../localcommands.tex}{
   \addbibresource{../localbibliography.bib}
   \usepackage{tabularx,multicol}
\usepackage[hyphens]{url}
\urlstyle{same}

\usepackage{langsci-optional}
\usepackage{langsci-lgr}
\usepackage{langsci-gb4e}

\usepackage{soul}
% \usepackage{booktabs}
% \usepackage{geometry}
\usepackage{multirow}
% \usepackage{makecell} %safely breaks a line inside a table cell
% \usepackage{tablefootnote} %footnotes in table captions
\usepackage{enumerate}
\usepackage{tikz}
\usepackage{array}
\usepackage[shortlabels]{enumitem}

%for having several figures on one line, turning words in 90°angles
% \usepackage{graphicx}
\usepackage{subcaption}

% \usepackage{caption} %for making the caption line longer in concrete cases

% \usepackage{rotating} %rotate a whole page; for large tables

\usepackage{fontspec}
\newfontfamily\charis{Charis SIL} %for being able to use the symbol ‿ and a better command of some of the boldfaced diacritics in 01
\newfontfamily\japanesefont{Noto Serif CJK JP} % for Japanese script
% \usepackage{tipa}

% \usepackage{needspace} %to keep exs on one page

\usepackage{xcolor}

\usepackage{pgfplots}
\pgfplotsset{compat=1.18}

%%%%%   AVMs for 02	%%%%%
\usepackage{langsci-avm}
\avmsetup{switch = ?}


\usepackage{langsci-branding}

   %\newcommand*{\orcid}{}


\makeatletter
\let\thetitle\@title
\let\theauthor\@author
\makeatother

\newcommand{\togglepaper}[1][0]{
   \bibliography{../localbibliography}
   \papernote{\scriptsize\normalfont
     \theauthor.
     \titleTemp.
     To appear in:
     E. Di Tor \& Herr Rausgeberin (ed.).
     Booktitle in localcommands.tex.
     Berlin: Language Science Press. [preliminary page numbering]
   }
   \pagenumbering{roman}
   \setcounter{chapter}{#1}
   \addtocounter{chapter}{-1}
}

\newbool{bookcompile}
\booltrue{bookcompile}
\newcommand{\bookorchapter}[2]{\ifbool{bookcompile}{#1}{#2}}



\renewcommand{\thempfootnote}{\fnsymbol{mpfootnote}} % use * for footnotes in float env

% \newcolumntype{L}[1]{>{\raggedright\arraybackslash}p{#1}} %left-aligned (non-justified) text in tables
% \newcolumntype{R}[1]{>{\raggedleft\arraybackslash}p{#1}} % right-aligned (non-justified) text in tables

\definecolor{customgreen}{RGB}{27,158,119}
\definecolor{custompink}{RGB}{231,42,138}
\definecolor{customblue}{RGB}{117,112,179}
\definecolor{customlightblue}{RGB}{143,170,220}
\definecolor{customdarkblue}{RGB}{68,114,196}
\definecolor{customyellow}{RGB}{225,192,0}
\definecolor{customorange}{RGB}{233,137,21}
\definecolor{customgray}{RGB}{229,229,229}


% Left-aligned indexes
% % \makeatletter
% % \patchcmd{\theindex}{\twocolumn}{\raggedright\twocolumn}{}{}
% % \makeatother

%to make Japanese scripts in the bibliography smaller
\newcommand{\japbib}[1]{{\japanesefont\small #1}}

%%%%%%%%%%%%%%%%%%%	MACROS for 02	%%%%%%%%%%%%%%%%

%%% Formatting	%%%%%

\newcommand{\textex}[1]{\textit{#1}} 		 %for formatting linguistic examples in-text%
%command-shift X

\newcommand{\lxm}[1]{\textsc{#1}}  		%for formatting names of lexemes
%command-option-shift L

%%%%%%%%% Abbreviations	%%%%%%

\newcommand{\featname}[1]{\textsc{#1}}
\newcommand{\featspec}[2]{\featname{#1}:{#2}}
\newcommand{\feat}[2]{\textsc{#1}:{#2}}  				% \feat{FEATURE}{value}  sc's
\newcommand{\featnameonly}[1]{\textsc{#1}} 			% \featnameonly{featurename}  sc's

\newcommand{\pr}{$^{\prime}$}

\newcommand{\Corr}{\textit{Corr}}
\newcommand{\propm}{\textit{pm}}

\newcommand{\PC}{Π$_{C}$}

\newcommand{\PF}{Π$_{F}$}
\newcommand{\PR}{Π$_{R}$}

\newcommand{\lex}{$\mathcal{L}$}

   %% hyphenation points for line breaks
%% Normally, automatic hyphenation in LaTeX is very good
%% If a word is mis-hyphenated, add it to this file
%%
%% add information to TeX file before \begin{document} with:
%% %% hyphenation points for line breaks
%% Normally, automatic hyphenation in LaTeX is very good
%% If a word is mis-hyphenated, add it to this file
%%
%% add information to TeX file before \begin{document} with:
%% %% hyphenation points for line breaks
%% Normally, automatic hyphenation in LaTeX is very good
%% If a word is mis-hyphenated, add it to this file
%%
%% add information to TeX file before \begin{document} with:
%% \include{localhyphenation}
\hyphenation{
    Al-ex-an-der
    Ber-ber
    chan-ges
    cir-cum-stan-ces
    coin-age
    Fraj-zyngier
    ita-liano
    Ja-pa-nese
    Kö-nig
    Krzy-żanowski
    li-te-ra-tur-no-go
    mi-zen-kei
    par-a-digm
    phe-nom-e-non
    ren-tai-shi
    Slo-vene
    Spen-cer
    Stoj-kov
    Ver-bi-ca-re-se
    wide-spread
}

\hyphenation{
    Al-ex-an-der
    Ber-ber
    chan-ges
    cir-cum-stan-ces
    coin-age
    Fraj-zyngier
    ita-liano
    Ja-pa-nese
    Kö-nig
    Krzy-żanowski
    li-te-ra-tur-no-go
    mi-zen-kei
    par-a-digm
    phe-nom-e-non
    ren-tai-shi
    Slo-vene
    Spen-cer
    Stoj-kov
    Ver-bi-ca-re-se
    wide-spread
}

\hyphenation{
    Al-ex-an-der
    Ber-ber
    chan-ges
    cir-cum-stan-ces
    coin-age
    Fraj-zyngier
    ita-liano
    Ja-pa-nese
    Kö-nig
    Krzy-żanowski
    li-te-ra-tur-no-go
    mi-zen-kei
    par-a-digm
    phe-nom-e-non
    ren-tai-shi
    Slo-vene
    Spen-cer
    Stoj-kov
    Ver-bi-ca-re-se
    wide-spread
}

   \boolfalse{bookcompile}
   \togglepaper[23]%%chapternumber
}{}

\begin{document}
\maketitle


\il{Slavic!South|(}
\section{Introduction} 
The loss of inflectional\is{inflection} case marking\is{case!marking} is an ongoing process in the dialect continuum\is{dialect!continuum} that stretches across \ili{Serbian} and \ili{Bulgarian} linguistic territory. This has led to substantial \isi{variation} in the \isi{dialect} case systems\is{case!system} of the region. We observe the first signs of this historical change in systems with rich \isi{case} distinctions, where case marking\is{case!marking} is obligatory on all classes of nominals\is{nominal}, as in the western part of the continuum\is{dialect!continuum}. Further east, we find dramatically reduced case paradigms\is{case!paradigm}, with just a few distinct \isi{case} forms, attested with only part of the lexicon, while the rest of the lexicon has lost the ability to inflect for \isi{case}. At the same time, in these dialects\is{dialect} we find competing rules of case marking\is{case!marking}, which appear to reflect different historical periods in the breakdown of the case system\is{case!system}. The result is either competition between more or less distinct \isi{inflection} (towards the west, in \ili{Serbian} dialects\is{dialect}), or competition between \isi{inflection} and uninflectedness (towards the east, in \ili{Bulgarian} dialects\is{dialect}). We investigate factors which underlie this competition and show how they vary at different stages of case decline\is{case!decline}. This provides us with an insight into the historical mechanisms underlying case loss\is{case!loss of case}, and to the conditions that may foster this change.

While our analysis has confirmed the relevance of the factors previously proposed to explain variable case marking\is{case!marking} synchronically\is{synchrony} and diachronically\is{diachrony} (e. g. \citealt{moravcsik1978, bossong1991, konig2008, dalrymplenikolaeva2011, cristofaro2013, kurumadajaeger2015, witzlackmakarevichserzant2018}), the data obtained from \ili{Serbian} and \ili{Bulgarian} dialects\is{dialect} shed new light on these factors and show that they do not in fact operate in the way that has previously been suggested.

First, extensive typological data show that explicit case marking\is{case!marking} is used in environments where it is essential to avoid ambiguity (e.g. between subjects\is{subject} and objects\is{object}) and is omitted where syntactic roles\is{syntactic!role} are transparent \citep{comrie1977}. For example, it has been shown that the omission of morphological case marking\is{case!marking} goes hand-in-hand with the use of adpositions\is{adposition}, which to a certain extent duplicate the information conveyed by \isi{case} \isi{inflection}. This has been shown to play a role in diachronic\is{diachrony} processes affecting case systems\is{case!system}: the information redundancy of \isi{case} \isi{inflection} in adpositional constructions\is{adpositional construction} may make these constructions more prone to the loss of morphological case\is{case!loss of case} than those without adpositions, where \isi{case} \isi{inflection} provides critical information about syntactic roles\is{syntactic!role} \citep{hewsonbubenik2006}. In the dialects\is{dialect} of East Serbia, while the use of adpositions\is{adposition} (prepositions\is{preposition}) correlates with a decline in case distinctions\is{case!decline}, the dialects\is{dialect} display a pattern which is the opposite of what is predicted by the above principle: the original case marking\is{case!marking} is better preserved precisely when governed by a \isi{preposition}, and more prone to loss elsewhere.

{\sloppy Second, pragmatic requirements, in particular \isi{discourse} salience, have been proven to be another important factor in variable case marking\is{case!marking} cross-linguistically. To mention just two examples, \citet{mcgregor2018} and \citet{montaut2018} present somewhat similar situations in unrelated languages (Southern African Khoisan\il{Khoisan!Southern African} and \ili{Dravidian} respectively) where the rise of \isi{object} case marking\is{case!marking} led to pragmatically motivated choices between overt case marking\is{case!marking} (\isi{accusative}) and zero marking: prominent objects\is{object} receive overt case marking\is{case!marking} while backgrounded objects\is{object} do not. In \ili{Bulgarian} border-region dialects\is{dialect}, where just two cases\is{case} have been preserved (\isi{nominative} and \isi{accusative}), explicit case marking\is{case!marking} on non-subject\is{subject!non-subject} nouns\is{noun} (\isi{accusative}) is also conditioned by pragmatic factors\is{pragmatic!factor}. However, it is not associated with pragmatic salience\is{pragmatic!salience}, but rather the opposite: explicit (\isi{accusative}) marking on non-subjects\is{subject!non-subject} is more frequent in pragmatically less prominent parts of an utterance (e.g. containing background information). Pragmatic prominence instead tends to trigger \isi{nominative} marking on non-subject\is{subject!non-subject} nouns\is{noun}, which makes them morphologically indistinguishable from subjects\is{subject}.\par}

We present two \isi{case} studies that show how different factors shape the transition from a larger to a smaller case system\is{case!system}, and further to uninflectedness. After the essential background information in Section \ref{05sec2}, we consider the role of structural factors, which are dominant in \ili{Serbian} dialects\is{dialect} (Section \ref{05sec3}), and analyse the impact of pragmatic conditions\is{pragmatic!condition}, which are dominant in the \ili{Bulgarian} dialects\is{dialect} under study (Section \ref{05sec4}). We suggest that this shift in prominence is due to the diachronic\is{diachrony} stage each variety has reached: structural factors are prominent at earlier stages of case decline\is{case!decline} and pragmatic factors\is{pragmatic!factor} at an advanced stage.

\section{Background} \label{05sec2}

\largerpage[-1.5]

South Slavic varieties inherited from Proto-Slavic\il{Slavic!Proto-Slavic} a case system\is{case!system} with six morphological cases\is{case} (plus the vocative whose status as \isi{case} is debatable\footnote{Vocative forms continue to be widely used in contemporary \ili{Serbian} and \ili{Bulgarian}, but do not express grammatical relationships \citep[8]{blake2004} and do not interact with the types of configurations we have investigated.}) and six major inflection classes\is{inflectional class} (\citealt[147--160]{mircev1958}; \citealt[127--178]{slavova2017}; \citealt[115]{popovic1955}). The further history of \isi{case} largely correlates with the division of the South Slavic territory into eastern and western zones. In the western zone (contemporary \ili{Serbian}, \ili{Croatian}, \ili{Bosnian}, \ili{Montenegrin}, \ili{Slovenian}) the original Proto-Slavic\il{Slavic!Proto-Slavic} case system\is{case!system} has been largely preserved. In the eastern zone (contemporary \ili{Bulgarian} and \ili{Macedonian}), most varieties have seen the reduction and loss of \isi{nominal} case marking\is{case!loss of case} and the rise of analytical constructions\is{analytical construction} (e.g. prepositional constructions\is{prepositional!construction}), as can be seen by contrasting the data from standard \ili{Bulgarian} and standard \ili{Serbian} in Table~\ref{05tab1}.

\begin{table}[ht]
	\small
	\centering
	\setlength{\tabcolsep}{4pt}
	\begin{tabularx}{\textwidth}{QQl}%L{4.1cm}L{3.4cm}L{4.3cm}
		\lsptoprule
		\textbf{Standard Serbian} \newline inflected forms & \textbf{Standard Bulgarian} \newline prepositional phrases\is{prepositional!phrase} with an uninflected form & \textbf{Translation}\\
		\midrule
		Ovo je \textbf{Kipar}. (\textsc{nom}) & Tova e \textbf{Kipâr}. & ‘This is Cyprus.’ \\
		Idu na \textbf{Kipar}. (\textsc{acc}) & Otivat na \textbf{Kipâr}. & ‘They go to Cyprus.’ \\
		Pomažu \textbf{Kipru}. (\textsc{dat}) & Pomagat na \textbf{Kipâr}. & ‘They help Cyprus.’ \\
		Stanovništvo \textbf{Kipra}. (\textsc{gen}) & Naselenieto na \textbf{Kipâr}. & ‘The population of Cyprus.’ \\
		Žive na \textbf{Kipru}. (\textsc{loc}) & Živejat na \textbf{Kipâr}. & ‘They live in Cyprus.’ \\
		Upravljaju \textbf{Kiprom}. (\textsc{ins}) & Upravljavat \textbf{Kipâr}. & ‘They govern Cyprus.’ \\
		\lspbottomrule
	\end{tabularx}
	\caption{Nominal\is{nominal} \isi{inflection} in \ili{Serbian} and corresponding forms in \mbox{\ili{Bulgarian}}\protect\footnotemark}
	\label{05tab1}
\end{table}

% \needspace{3\baselineskip}

\footnotetext{\ili{Bulgarian} demonstrates an asymmetry between \isi{nominal} and pronominal systems with respect to case marking\is{case!marking}. Nouns\is{noun} and adjectives\is{adjective} do not inflect for \isi{case} (in standard varieties and in a significant number of dialects\is{dialect}). Personal pronouns distinguish full and clitic forms (the choice being syntactically determined). Full personal pronouns retain a two-\isi{case} distinction, \isi{nominative} vs. \isi{accusative}, the latter being used in all non-subject\is{subject!non-subject} syntactic roles\is{syntactic!role}: \emph{\textbf{tja}} (\textsc{3sg.f.nom}) \emph{običa} (`she loves') vs. \emph{običam \textbf{neja}} (\textsc{3sg.f.acc}) (`I love her') vs. \emph{kazax na \textbf{neja}} (\textsc{3sg.f.acc}) (`I told her', literally `to her'). Clitic personal pronouns are used in non-subject\is{subject!non-subject} roles and distinguish between \isi{accusative} and \isi{dative}: \emph{običam \textbf{ja}} (\textsc{3sg.f.acc}) vs. \emph{kazax \textbf{ì}} (\textsc{3sg.f.dat}) (`I told her').}

In between these two types there is a substantial transitional \isi{dialect} zone. A characteristic feature of these transitional dialects\is{dialect} is the competition between these two strategies, as seen both in \ili{Serbian} (\ref{05ex1}) and \ili{Bulgarian} dialects\is{dialect} (\ref{05ex2}).

\begin{exe}
	\ex \label{05ex1}
	\gll Zove ga otac telefon-om / na telefon. \\
	call[\textsc{prs.3sg}] 3\textsc{sg.m.acc} father[\textsc{nom}] phone-\textsc{ins.sg} / on phone[\textsc{acc}](\textsc{=nom}) \\
	\trans `The father calls him by phone.' \hfill \citep[183, 202]{miloradovic2003}
	
	\ex \label{05ex2}
	\gll Petr-a / Petâr go njama. \\
	Peter-\textsc{sg.acc} / Peter[\textsc{sg.nom}] 3\textsc{sg.m.acc} \textsc{exist.neg} \\
	\trans `Peter is absent.' \hfill \citep[163]{stojkov1981}
\end{exe}

% \largerpage[2]

The dialects\is{dialect} in this transitional area correspond to different postulated chronological stages of the decline of the Proto-Slavic\il{Slavic!Proto-Slavic} case system\is{case!system}. A group of dialects\is{dialect} in the west of the transitional area preserve all six original Proto-Slavic\il{Slavic!Proto-Slavic} cases\is{case} (plus vocative), but speakers regularly replace three of them with the \isi{accusative} form, which serves as a general oblique\is{oblique!general} \isi{case} (Figure \ref{05tab2}, Transitional System 1). Thus, the original six-\isi{case} system competes in these dialects\is{dialect} with an innovative three-\isi{case} system, consisting of \isi{nominative}, \isi{dative}, and \isi{accusative} only (\citealt{miloradovic2003, simic1972}). The further east one goes, the greater are the reductions. The \isi{dative} \isi{case} is the next to be replaced (Transitional System 2). Here, the more conservative three-\isi{case} system competes with a two-\isi{case} system, consisting of \isi{nominative} and \isi{accusative} (\citealt{belic1905, ivic1956}). Finally, in the east of the transitional area, the \isi{accusative} itself starts to yield to a single uninflected form (Transitional System 3). Here, the two-\isi{case} system coexists with a system without case\is{case!system} \citep{stojkov1975, stojkov1981}. Thus, three structural \isi{dialect} types can be distinguished within the transitional area, according to the size of the case paradigm\is{case!paradigm}: six-, three- and two-\isi{case} systems, as summarised in Figure \ref{05tab2}.\footnote{In distinguishing between the largely syncretic\is{syncretism} \isi{dative} and \isi{locative} cases\is{case} in \ili{Serbian}, we follow the Slavic tradition of \isi{case} description (\citealt{piperklajn2013, miloradovic2003}). The non-syncretic\is{syncretism} \isi{locative} endings were used at least until the end of the 15\textsuperscript{th} century \citep{belic1999}, while the loss of case\is{case!loss of case} started sometime before the 12\textsuperscript{th} century \citep{belic1905}. It follows, then, that when the process of the case loss\is{case!loss of case} started, the \isi{dative} and \isi{locative} were still morphologically distinct cases\is{case}. An additional argument in favour of differentiating between the two cases\is{case} comes from the fact that it facilitates the description of the process of case loss\is{case!loss of case}. The original \isi{locative} \isi{case} is frequently replaced by the general oblique\is{oblique!general} in the whole transitional area, including the western region where it is the only \isi{case} to be replaced. The original \isi{dative} \isi{case}, on the contrary, is preserved in most \ili{Serbian} dialects\is{dialect} with the exception of the eastern zone where the only remaining cases\is{case} are \isi{nominative} and \isi{accusative} \citep{sobolev1991}.\label{05ftn3}}

\begin{figure}
	\centering
	\includegraphics[width=\textwidth]{figures/05tab2.pdf}
	\caption{Case systems\is{case!system} in the \ili{Serbian}-\ili{Bulgarian} dialect continuum\is{dialect!continuum}}
	\label{05tab2}
\end{figure}

The chronology that this transition reflects has been reconstructed primarily on the basis of liturgical texts produced within various regional written traditions in the South Slavic territory from the 10\textsuperscript{th} century onwards (\citealt{tsonev1984, duridanov1956, duridanov1958, rusek1964, cesko1970, belic1905}) and, for later periods, on the basis of clerical and household documents \citep{bernstejn1948}.

It is believed that the decay of the case system\is{case!system} originated in the South Slavic eastern zone (\ili{Bulgarian} and \ili{Macedonian}) around the 10\textsuperscript{th} century \citep[25]{duridanov1958}. By the end of the 14\textsuperscript{th} century morphological case marking\is{case!marking} on nominals\is{nominal} was completely lost in at least some of the dialects\is{dialect} spoken in the eastern zone (\ili{Bulgarian} border-region dialects\is{dialect}) \citep[11]{cesko1970}. Even earlier, by the end of the 13\textsuperscript{th} century, case loss\is{case!loss of case} was at an advanced stage in a number of \ili{Macedonian} dialects\is{dialect} \citep{rusek1964}. In \ili{Serbian}, the process of case decline\is{case!decline} started sometime before the 12\textsuperscript{th} century and affected the South-East dialects\is{dialect} \citep{belic1905}. It did not, however, lead to the complete loss of \isi{case} distinctions: rather it resulted in different types of reduced systems.


Historical evidence points to the following processes involved in this change: the spread of prepositional constructions\is{prepositional!construction} replacing bare noun phrases\is{noun!phrase}; the redistribution of functions between indirect cases\is{case!indirect} and displacement of less productive indirect cases\is{case!indirect}; the spread of \isi{accusative} forms in the function of an oblique case\is{case!oblique} generalising across original indirect cases\is{case!indirect};\footnote{The chronology of the displacement of individual indirect cases\is{case!indirect} varies. The most stable, as dialectological\is{dialect} data of the \ili{Serbian} language show (\citealt{belic1905, miloradovic2003}) is the \isi{dative} \isi{case}. The historical data from \ili{Bulgarian} also testify to this: the \isi{dative} forms of nouns\is{noun} were found in \ili{Bulgarian} texts of the 17\textsuperscript{th}, 18\textsuperscript{th} and even 19\textsuperscript{th} centuries, while the \isi{genitive}, \isi{instrumental} and local cases\is{case!loss of case} had already fallen out of use in \ili{Bulgarian} dialects\is{dialect} in the 13\textsuperscript{th}-14\textsuperscript{th} centuries (\citealt[255--260]{mircev1958}; \citealt[301--310]{cesko1970}).} and the spread of uninflected forms replacing this oblique case\is{case!oblique}.

\section{Constructional \isi{variation}: the loss of the \isi{genitive} in \ili{Serbian} dialects\is{dialect}} \label{05sec3}

\largerpage[1]

The loss of case\is{case!loss of case} in European languages has been reported to go hand in hand with the increased use of prepositions\is{preposition} (\citealt{blake2004, hewsonbubenik2006, kulikov2006}). \citet[178]{blake2004} illustrates this with an Old English\il{English!Old} version of Luke 15:15 (\ref{05ex3a}) and Wyclif's fourteenth-century translation (\ref{05ex3b}) (glossing and translation are as in \citealt{blake2004}):

\begin{exe}
	\ex
	\begin{xlist}
		\ex \label{05ex3a}
		\gll He folgode an-um burg-sitt-end-um menn thaes rices. \\
		he followed one-\textsc{dat} town-dwell-ing-\textsc{dat} man.\textsc{dat} that.\textsc{gen} land.\textsc{gen} \\
		\ex \label{05ex3b}
		He clevede [=attached himself] to oon of the citizens of that contré. \\
		\trans ‘He hired himself out to one of the citizens of that country.’
	\end{xlist}
\end{exe}

In (\ref{05ex3a}), the directional (`to the citizen') and the \isi{possessive} (`of that country') relationships are expressed via \isi{inflection}, while in (\ref{05ex3b}), they are expressed by prepositional constructions\is{prepositional!construction}. The relationship between the two processes, however, is not clear. Did the spread of prepositional constructions\is{prepositional!construction} cause case\is{case!loss of case} to disappear as it became semantically redundant (as claimed in \citealt{hjelmslev1935, hewsonbubenik2006})? Or was it, by contrast, a repair strategy necessary after the loss of explicit case\is{case!loss of case} marking, as proposed, for example, in \citet{skautrup1944}?

The \isi{variation} found in \ili{Serbian} dialects\is{dialect} allows us to investigate the interplay of inflectional\is{inflection} case marking\is{case!marking} and prepositional constructions\is{prepositional!construction} using synchronic\is{synchrony} data. We focus here on the \isi{genitive}, which is in the process of being replaced by the \isi{accusative}. Specifically, we ask: do prepositional constructions\is{prepositional!construction} particularly favour this replacement? This is what an \isi{informativeness} principle would predict: since the \isi{preposition} provides semantic context, \isi{genitive} case marking\is{case!marking} is redundant and the more general \isi{accusative} \isi{case} can be used without any loss of information (cf. \citealt{comrie1977, comrie1989}). Our data, however, show the opposite tendency: prepositional constructions\is{prepositional!construction} favour the \emph{retention} of the \isi{genitive}, while its replacement by the \isi{accusative} is more advanced in non-prepositional contexts\is{prepositional!construction}. That means that in \ili{Serbian} dialects\is{dialect}, some principle other than \isi{informativeness} is at play. We suggest below that this principle could be understood as a speaker-oriented strategy of grammatical change.

In Section \ref{05sec3.1}, we introduce the data used in this \isi{case} study. Section \ref{05sec3.2} outlines the use of the \isi{genitive} forms in \ili{Serbian} and discusses the results of previous research on the \isi{variation} between \isi{genitive} and \isi{accusative} in South-East \ili{Serbian} dialects\is{dialect}. Sections \ref{05sec3.3} and \ref{05sec3.4} present the patterns of \isi{variation} found in our data; Section \ref{05sec3.5} discusses the theoretical implications of our findings.

\subsection{Data} \label{05sec3.1}

The data for this study come from sociolinguistic interviews with speakers of \ili{Serbian} dialects\is{dialect} located in or near Kosovo. These come from two sources: the archive of the Institute for Balkan Studies and our own fieldwork. The archive provided a subset of recordings collected in the project ``Research of Slavic Vernaculars in Kosovo and Metohija'' (2002 -- 2003) led by the Institute for the \ili{Serbian} Language of SASA and financed by UNESCO. Our data come from fieldwork in the municipality of Brus, Central Serbia, conducted by Maria Kyuseva within the project ``Declining case: inflectional\is{inflection} loss in progress''\is{case!decline} in 2022.

These data were collected in a uniform fashion that conforms to the \ili{Serbian} dialectological\is{dialect} tradition. The participants were native speakers of local dialects\is{dialect}, aged between 70 and 81 years, who had been long-term residents of the region and had engaged in traditional activities throughout their lives, such as cattle breeding or agriculture. This sampling ensured the minimal interference of the high-prestige standardised variety of \ili{Serbian}. The topics\is{topic} of the interviews included history, tradition, culture, crafts, cuisine, everyday life, and biographical stories.

The total corpus compiled from these two sources contains over 274,800 tokens, 30 hours and 36 minutes of recordings. Overall, twenty-three speakers from eleven villages were recorded. Figure \ref{05fig1} shows the locations of the villages, which are colour-coded according to \isi{dialect}. The Zeta-South Sandžak \isi{dialect} is marked in green, the Kosovo-Resava \isi{dialect} is marked in pink, and the Prizren-South Morava \isi{dialect} is marked in purple.

\begin{figure}
	\includegraphics[width=0.9\textwidth]{figures/05fig1.pdf}
	\caption{
		Locations of the interviews in Serbia
		\newline
		\newline
		\textcolor{customgreen}{\Large ●} Zeta–South Sandžak \quad
		\textcolor{custompink}{\Large ●} Kosovo–Resava \quad
		\textcolor{customblue}{\Large ●} Prizren–Timok
	}
	\label{05fig1}
\end{figure}

With respect to case marking\is{case!marking}, these three dialects\is{dialect} have the same system (Transitional System 1 in Figure \ref{05tab2}): all six original cases\is{case} are preserved to various degrees, but three of them (\isi{genitive}, \isi{locative}, and \isi{instrumental}) can be replaced by \isi{accusative} forms. The dialects\is{dialect} differ in the frequency of such replacements. Zeta-South Sandžak is the most conservative: it has a stable six-\isi{case} system, and the three peripheral cases\is{case} are only occasionally replaced by the \isi{accusative}. Kosovo-Resava represents a more advanced stage of case loss\is{case!loss of case}: the \isi{genitive}, \isi{locative}, and \isi{instrumental} cases\is{case} frequently give way to the \isi{accusative}. Finally, Prizren-Timok shows the most advanced stage: the \isi{locative} and \isi{instrumental} are virtually non-existent (occurring with only a handful of lexemes), and the use of the \isi{genitive} is highly restricted. We compare the usage of \isi{genitive} forms in these three dialects\is{dialect} and project the synchronic\is{synchrony} tendencies onto the diachronic\is{diachrony} plane. This allows us to draw conclusions regarding the historical dynamics involved in the decline of this case\is{case!decline}.

\is{genitive|(}
\subsection{Genitive \isi{case} in \ili{Serbian}} \label{05sec3.2}

The \ili{Serbian} genitive is highly polysemous. Its typical function is to mark a relationship between two entities, such as possessor-possessee, part-whole, and source-goal. Additionally, it can denote time, reason, goal, and manner. From the point of view of syntax, it has a wide distribution. It can be used in adnominal\is{adnominal construction}, adjectival\is{adjectival construction}, and verbal constructions\is{verbal construction}, as well as in prepositional constructions\is{prepositional!construction} and constructions\is{construction} with a \isi{quantifier}. Out of all the cases\is{case}, the genitive occurs with the widest inventory of prepositions\is{preposition} \citep{felesko1955}.

Various studies of individual dialects\is{dialect} of South-East Serbia have already established the heterogeneous behaviour of the genitive in different constructions\is{construction}. For example, \citet{miloradovic2003} analyses the competition between the genitive and \isi{accusative} in the Kosovo-Resava \isi{dialect} of the Paraćin region in Central Serbia. She observes that while the genitive is dominant in prepositional constructions\is{prepositional!construction}, its use is highly limited elsewhere. In verbal constructions\is{verbal construction}, the genitive is replaced by the \isi{accusative} (\ref{05ex4a}), in adnominal constructions\is{adnominal construction} -- by a prepositional phrase\is{prepositional!phrase} (\ref{05ex4b}), and in adjectival constructions\is{adjectival construction}, it can be replaced either by the \isi{accusative} or by a prepositional phrase\is{prepositional!phrase} (\ref{05ex4c}).

\begin{exe}
	\ex
	\begin{xlist}
		\ex \label{05ex4a}
		\gll Nije bi-l-o \textbf{vod-e}. $\rightarrow$ Nije bi-l-o \textbf{vod-u}. \\
		\textsc{aux.neg.3sg} be-\textsc{ptcp-n.sg} water-\textsc{gen.sg} \phantom{} \textsc{aux.neg.3sg} be-\textsc{ptcp-n.sg} water-\textsc{acc.sg} \\
		\trans ‘There was no water.’
		
		\ex \label{05ex4b}
		\gll kraj \textbf{mačug-e} $\rightarrow$ kraj \textbf{od} \textbf{mačug-e} / \textbf{od} \textbf{mačugu} \\
		end[\textsc{nom.sg}] cane-\textsc{gen.sg} \phantom{} end[\textsc{nom.sg}] of cane-\textsc{gen.sg} / of cane-\textsc{acc.sg} \\
		\trans ‘cane end’
		
		\ex \label{05ex4c}
		\gll Pun \textbf{vod-e}. $\rightarrow$ Pun \textbf{vod-u}. \\
		full[\textsc{m.nom.sg}] water-\textsc{gen.sg} \phantom{} full[\textsc{m.nom.sg}] water-\textsc{acc.sg} \\
		\gll \phantom{full.m.nom.sg water-gen.sg $\rightarrow$} Pun \textbf{s} \textbf{vod-u}. \\
		\phantom{full.[m.nom.sg] water-gen.sg $\rightarrow$} full[\textsc{m.nom.sg}] with water-\textsc{acc.sg} \\
		\trans ‘[It is] full of water.’
	\end{xlist}
\end{exe}

\citet{stankovic2022} analyses the speech of the younger population in the town of Vranje in South Serbia, where the Prizren-South Morava \isi{dialect} is spoken. This \isi{dialect} has a three-\isi{case} system consisting of the \isi{nominative}, \isi{accusative}, and \isi{dative}. The speech of the younger generation, however, is influenced by the standard \isi{dialect}. Therefore, it is susceptible to the penetration of other cases\is{case} found in the standard, especially the genitive. Stanković observes that the possibility of using the genitive depends on the \isi{construction}. Focusing on the non-prepositional\is{preposition} uses, she finds that adnominal constructions\is{adnominal construction} and constructions\is{construction} with a \isi{numeral} permit the occasional appearance of the genitive \isi{case} (\ref{05ex5a}-\ref{05ex5b}), while constructions\is{construction} where the complement depends on a \isi{verb} or a predicative \isi{adjective} do not (\ref{05ex5c}-\ref{05ex5d}):

\begin{exe}
	\ex
	\begin{xlist}
		
		\ex \label{05ex5a}
		\gll udic-a za \textbf{pecanj-e} \textbf{rib-a} \\
		rod-\textsc{nom.sg} for catching-\textsc{acc.sg} fish-\textsc{gen.pl} \\
		\trans ‘fishing rod’
		
		\ex \label{05ex5b}
		\gll Ima-m \textbf{pet} \textbf{različit-ih} \textbf{igar-a}. \\
		have-\textsc{prs.1sg} five different-\textsc{gen.pl} game-\textsc{gen.pl} \\
		\trans ‘I have five different games.’
		
		\ex \label{05ex5c}
		\gll Nije \textbf{je-o} \textbf{čokolad-u} / *čokolad-e. \\
		\textsc{aux.neg.3sg} eat-\textsc{ptcp.m.sg} chocolate-\textsc{acc.sg} / *chocolate-\textsc{gen.sg} \\
		\trans ‘He didn’t eat the chocolate.’
		
		\ex \label{05ex5d}
		\gll Dvorišt-e \textbf{pun-o} \textbf{s} \textbf{igračk-e} / *pun-o igračak-a. \\
		yard-\textsc{nom.sg} full-\textsc{n.nom.sg} \textbf{with} toy-\textsc{acc.pl} / *full-\textsc{n.nom.sg} toy-\textsc{gen.pl} \\
		\trans ‘The yard is full of toys.’
	\end{xlist}
\end{exe}

In our study, we bring these individual observations together by comparing the use of the genitive \isi{case} forms in different constructions\is{construction} across the three South-East \ili{Serbian} dialects\is{dialect}: Zeta-South Sandžak, Kosovo-Resava, and Prizren-South Morava.

\is{syntactic!factor|(}
\subsection{Constructions\is{construction} with the genitive in South-East \ili{Serbian}} \label{05sec3.3}

We focus on three constructions\is{construction} with the genitive \isi{case}: prepositional constructions\is{prepositional!construction}, constructions\is{construction} with a \isi{quantifier}, and adnominal constructions\is{adnominal construction}. In South-East \ili{Serbian} dialects\is{dialect}, these constructions\is{construction} allow both the original genitive and the innovative \isi{accusative} forms:

\begin{exe}
% \needspace{4\baselineskip}
\ex Prepositional construction\is{prepositional!construction}
\begin{xlist}
	\ex 
	\gll Napravljen-i \textbf{od} \textbf{drv-a}. \\
	made-\textsc{m.nom.pl} from wood-\textsc{gen.sg} \\
	\trans ‘[They are] made from wood.’
	
	\ex 
	\gll Ima i \textbf{od} \textbf{drv-o}. \\
	\textsc{exist.prs} and from wood-\textsc{acc.sg} \\
	\trans ‘There is [some of it] from wood as well.’
\end{xlist}
\end{exe}

\begin{exe}
\ex Construction\is{construction} with a \isi{quantifier}
\begin{xlist}
	\ex 
	\gll Malo \textbf{trav-e} gi stavi-m. \\
	a.little grass-\textsc{gen.sg} 3\textsc{pl.dat} put-\textsc{prs.1sg} \\
	\trans ‘I put out a bit of grass for them.’
	
	\ex 
	\gll I ondak turi malo \textbf{trav-u}. \\
	and then put[\textsc{imp}] a.little grass-\textsc{acc.sg} \\
	\trans ‘And then put out a little bit of grass.’
\end{xlist}
\end{exe}

\begin{exe}
\ex Adnominal construction\is{adnominal construction}
\begin{xlist}
	\ex 
	\gll Tamo odnese-š čaš-u \textbf{vin-a}. \\
	there take.away-\textsc{prs.2sg} cup-\textsc{acc.sg} wine-\textsc{gen.sg} \\
	\trans ‘You take there a cup of wine.’
	
	\ex 
	\gll I uzme-š čaš-u \textbf{vin-o}. \\
	and take-\textsc{prs.2sg} cup-\textsc{acc.sg} wine-\textsc{acc.sg} \\
	\trans ‘And you take a cup of wine.’
\end{xlist}
\end{exe}

The overall number of examples in the dataset is 960. The prepositional construction\is{prepositional!construction} is represented by 676 examples, the adnominal construction\is{adnominal construction} -- by 169 examples, and the construction\is{construction} with a \isi{quantifier} -- by 115 examples. The ratio represents the occurrence of these constructions\is{construction} in spontaneous speech. For the prepositional construction\is{prepositional!construction}, we include only phrases with the \isi{preposition} \emph{od} `from, of', which is the most frequent \isi{preposition} that governs the genitive. Constructions\is{construction} with a \isi{quantifier} include the following quantifiers\is{quantifier}: \emph{mnogo} `a lot', \emph{malo} `a little', \emph{lek} `very little', \emph{pomalo} `a little bit', \emph{puno} `much', \emph{više} `more', and \emph{manje} `less'. Finally, adnominal constructions\is{adnominal construction} include a variety of governing nouns\is{noun}, the most frequent meanings in the sample are: quantity (a cup of tea, a bowl of rice), and part-whole (top of the table, edge of the village). True possession (a son of my brother), otherwise typical for the \isi{construction}, is rare in our examples. It is usually expressed in the dialects\is{dialect} by alternative means of expression, such as \isi{possessive} adjectives\is{adjective} and constructions\is{construction} with the \isi{dative} \isi{case}.

The bar chart (Figure \ref{05fig2}) shows the \isi{variation} between the original genitive and the innovative \isi{accusative} in the three dialects\is{dialect}.

\begin{figure}
	\centering
	\includegraphics[width=0.546\textwidth]{figures/05fig2a.pdf}
	\hfill
	\includegraphics[width=0.412\textwidth]{figures/05fig2b.pdf}
	\includegraphics[width=0.525\textwidth]{figures/05fig2c.pdf}
	\caption{Competing strategies of case marking\is{case!marking} in South-East \ili{Serbian} dialects\is{dialect}}
	\label{05fig2}
\end{figure}

As the figure illustrates, in Zeta-South Sandžak (the top left section of the bar chart), the most conservative \isi{dialect} in the sample, the genitive prevails in all three constructions\is{construction}. It appears in around 90\% of the examples in the prepositional\is{prepositional!construction} and adnominal constructions\is{adnominal construction} and is used somewhat less in the \isi{construction} with a \isi{quantifier}. As an ANOVA\footnote{The ANOVA analysis of the corpus data here and below was carried out by Dr Alexander Stewart (University of St Andrews, UK).} test shows, there is no significant effect of the \isi{construction} on the \isi{case} in this \isi{dialect}, with p=0.2429 and F=1.4284. The difference between the three constructions\is{construction} is more prominent in the Kosovo-Resava \isi{dialect}, which has a more advanced stage of case decline\is{case!decline}. Here, the genitive still dominates with 90\% of occurrences in the prepositional construction\is{prepositional!construction}, it is slightly less frequent in the adnominal construction\is{adnominal construction}, and its use falls dramatically in the \isi{construction} with a \isi{quantifier}, where it is no longer a dominant strategy. The ANOVA test shows a significant effect of the \isi{construction} on the \isi{case} in this \isi{dialect} with p=3.3092e-18 and F=44.1388. Finally, the most innovative Prizren-South Morava \isi{dialect} (the lower section of the bar chart) shows yet another picture. Here, the use of \isi{accusative} is the dominant strategy in both the adnominal construction\is{adnominal construction} (66\% of examples) and the \isi{construction} with a \isi{quantifier} (92\% of examples), and the only \isi{construction} where the genitive appears in a slight majority of examples (57\%) is the prepositional construction\is{prepositional!construction}. Again, ANOVA shows a significant effect of the \isi{construction} with p=1.1018e-06 and F=14.2548. (The null hypothesis for each of the three dialects\is{dialect} assumed no effect of the \isi{construction} on case selection\is{case!selection}.)

These data show that the use of the genitive versus the \isi{accusative} in the \ili{Serbian} dialects\is{dialect} under study is at least partly conditioned by the type of \isi{construction}. The exact conditions depend on the \isi{dialect}, or if we accept the diachronic\is{diachrony} interpretation of the data, on the stage of case decline\is{case!decline}. The general tendency is the following: constructions\is{construction} with a \isi{quantifier} are the most susceptible to the penetration of the innovative \isi{accusative}, adnominal constructions\is{adnominal construction} come second, and prepositional constructions\is{prepositional!construction} preserve the original genitive the longest. This effect is clear when we look at the distribution of cases\is{case} in individual dialects\is{dialect}, but it can also be seen if we trace the linguistic behaviour of the constructions\is{construction} across dialects\is{dialect}. The use of the genitive case consistently declines\is{case!decline} in each \isi{construction} as we move across the transitional zone from west to east. For each \isi{construction}, we find a significant effect of the \isi{dialect}: with p=2.1082e-25 and F=61.8976 (prepositional construction\is{prepositional!construction}), p=9.1127e-10 and F=23.6043 (adnominal construction\is{adnominal construction}), and p=1.1889e-06 and F=15.4478 (\isi{construction} with a \isi{quantifier})\is{genitive|)}\is{syntactic!factor|)}.

\is{morphological!factor|(}
\subsection{Morphologically conditioned\is{morphological!condition} \isi{variation} in \ili{Serbian} dialects\is{dialect}} \label{05sec3.4}

Another factor that conditions \isi{variation} between the original \isi{genitive} and innovative \isi{accusative} \isi{case} is \isi{inflection} class. This plays a role in the whole dialectal continuum\is{dialect!continuum}, but in our data it is especially pronounced in the Prizren-South Morava \isi{dialect}.

In standard \ili{Serbian}, nouns\is{noun} cluster into four inflection classes\is{inflectional class} (Table \ref{05tab3}). The paradigms\is{paradigm} exhibit a high degree of \isi{syncretism} (see the comment in footnote \ref{05ftn3} concerning \isi{dative}-\isi{locative} syncretic\is{syncretism} forms).

\begin{table}[ht]
	\centering
	\caption{Inflection classes\is{inflectional class} for selected nouns\is{noun}}
	\label{05tab3}
	\begin{tabularx}{\textwidth}{lllll}
		\lsptoprule
		& Inflection & Inflection & Inflection & Inflection  \\
		& class I & class II & class III & class IV \\
		\midrule
		& prostor ‘space’ & kuća ‘house’ & stvar ‘thing’ & polje ‘field’ \\
		& (\textsc{masc}) & (\textsc{fem} / some \textsc{masc}) & (\textsc{fem}) & (\textsc{neut}) \\
		\midrule
		\textsc{sg.nom} & prostor & kuć-a & stvar & polj-e \\
		\textsc{pl.nom} & prostor-i & kuć-e & stvar-i & polj-a \\
		\textsc{sg.acc} & prostor & kuć-u & stvar & polj-e \\
		\textsc{pl.acc} & prostor-e & kuć-e & stvar-i & polj-a \\
		\textsc{sg.gen} & prostor-a & kuć-e & stvar-i & polj-a \\
		\textsc{pl.gen} & prostor-ā & kuć-ā & stvar-ī & polj-ā \\
		\textsc{sg.dat} & prostor-u & kuć-i & stvar-i & polj-u \\
		\textsc{pl.dat} & prostor-ima & kuć-ama & stvar-ima & polj-ima \\
		\textsc{sg.loc} & prostor-u & kuć-i & stvar-i & polj-u \\
		\textsc{pl.loc} & prostor-ima & kuć-ama & stvar-ima & polj-ima \\
		\textsc{sg.ins} & prostor-om & kuć-om & stvar-ju & polj-em \\
		\textsc{pl.ins} & prostor-ima & kuć-ama & stvar-ima & polj-ima \\
		\lspbottomrule
	\end{tabularx}
\end{table}

In South-East \ili{Serbian} dialects\is{dialect}, there is a tendency for nouns\is{noun} of inflection class III\is{inflectional class} to adopt the forms (and the \isi{gender}) of inflection class I\is{inflectional class}. For this reason, we do not have enough examples to be able to draw reliable generalizations about this class, and we exclude it from further analysis. The bar chart in Figure \ref{05fig3} shows the \isi{variation} between the \isi{genitive} and the \isi{accusative} for the remaining three inflection classes\is{inflectional class} in the Prizren-South Morava \isi{dialect}.

\begin{figure}
	\centering
	\includegraphics[width=0.6\textwidth]{figures/05fig3.pdf}
	\caption{Prizren-South Morava \isi{dialect}: original \isi{genitive} vs. innovative \isi{accusative} in different inflection classes\is{inflectional class}}
	\label{05fig3}
\end{figure}

The choice of \isi{case} is clearly conditioned by inflection class\is{inflectional class}: class II favours the \isi{genitive} form, while class I and class IV favour the \isi{accusative} form. This pattern can be illustrated with examples of the same meaning of the \isi{genitive} (material) being expressed differently depending on the inflection class\is{inflectional class}, see (\ref{05ex9a}-\ref{05ex9b}).

\begin{exe}
	\ex
	\begin{xlist}
		\ex \label{05ex9a}
		\gll Se ispred-e \textbf{od} \textbf{vun-e}. \\
		\textsc{refl} spin-\textsc{prs.3sg} from wool\textsc{(II)-gen.sg} \\
		\trans ‘It is spun from wool.’
		
		\ex \label{05ex9b}
		\gll Pravi-l-i \textbf{od} \textbf{drv-o}. \\
		make-\textsc{ptcp-m.pl} from wood\textsc{(IV)-acc.sg} \\
		\trans ‘They used to make it from wood.’
	\end{xlist}
\end{exe}


In (\ref{05ex9a}), the noun\is{noun} belongs to class II (\emph{vuna} `wool') and takes the \isi{genitive}; in (\ref{05ex9b}), the noun\is{noun} belongs to class IV (\emph{drvo} `wood') and takes the \isi{accusative}. This tendency to preserve the original \isi{case} form in class II is even more pronounced in \ili{Bulgarian} border-region dialects\is{dialect}, where case marking\is{case!marking} with class II occurs most frequently and is semantically unrestricted (see Section \ref{05sec4})\is{morphological!factor|)}.


\is{syntactic!factor|(}
\subsection{Discussion: structural factors in case decline\is{case!decline}} \label{05sec3.5}

These data illustrate both syntactic\is{syntactic!condition} and morphological conditions\is{morphological!condition} on \isi{variation} between the \isi{genitive} and \isi{accusative}. In terms of syntax, certain constructions\is{construction} favour the use of the \isi{accusative} over the \isi{genitive}, in particular prepositional phrases\is{prepositional!phrase}. In terms of morphology, certain inflection classes\is{inflectional class} are more likely to adopt the innovative \isi{accusative}, in particular class II.

The tendency of prepositional phrases\is{prepositional!phrase} to preserve the original case marking\is{case!marking} contradicts commonly held assumptions about the role played by prepositions\is{preposition} in case loss\is{case!loss of case} (\citealt{blake2004, hewsonbubenik2006}). Since \isi{genitive} case marking\is{case!marking} in these prepositional constructions\is{prepositional!construction} is semantically redundant, it could be omitted without incurring any ambiguity. One might assume then that prepositional constructions\is{prepositional!construction} would be the first to lose the \isi{genitive}. The fact that our data show the opposite means that the preservation of the \isi{genitive} is not driven by the need to disambiguate. We propose here a possible alternative motivation.

Following \citet{norde2002}, we distinguish two types of less-effort strategies for grammatical change: a speaker-oriented and a hearer-oriented strategy. These can be thought of as two opposing forces that shape any language change. The speaker ``wants'' to make less effort in encoding the message, and the hearer ``wants'' to make less effort in decoding the message. The \isi{informativeness} principle aligns with a hearer-oriented strategy. In this instance, the \isi{genitive} should first be replaced in prepositional constructions\is{prepositional!construction}, as it is not needed for the hearer to be able to correctly interpret the message, while it has a function in other constructions\is{construction}. А speaker-oriented strategy would be one that favours ease of production. The \isi{genitive} \isi{case} after the \isi{preposition} \emph{od} `from' is syntactically determined: \emph{od} governs the \isi{genitive}, creating a high degree of \isi{predictability}: every time the \isi{preposition} \emph{od} `from' is used, it can be expected to be followed by the \isi{genitive}. In adnominal constructions\is{adnominal construction} and constructions\is{construction} with a \isi{quantifier}, there is no such syntactic determination, in the sense that the nouns\is{noun}, such as `cup' or `edge', do not obligatorily govern a noun phrase\is{noun!phrase}, and quantifiers\is{quantifier} can similarly be used on their own. In this way, the syntactic information supports the \isi{case} morphology and contributes to the preservation of the original \isi{case} forms. Crucially, there is no gain for the hearer in preserving the \isi{case} forms in prepositional constructions\is{prepositional!construction} while omitting them in non-prepositional constructions. Rather, it is in the speaker's interest to keep an inflected form in a context that always requires it.

With these results in mind, we turn to the contrasting patterns of \ili{Bulgarian} border-region dialects\is{dialect}\is{syntactic!factor|)}.

\section{Competing case marking\is{case!marking} in \ili{Bulgarian} border-region dialects\is{dialect}: pragmatic factors\is{pragmatic!factor}} \label{05sec4}

In \ili{Bulgarian} border-region dialects\is{dialect}, case marking\is{case!marking} is lost in all inflection classes\is{inflectional class} of nouns\is{noun} except class II and, to some extent, class I (restricted to animate nouns\is{noun}). In these classes, singular nouns\is{noun} retain two \isi{case} forms, \isi{nominative} and \isi{accusative}; plural forms are not inflected for \isi{case}. The \isi{nominative} marks nouns\is{noun} in the \isi{subject} position and the \isi{accusative} marks nouns\is{noun} in non-subject\is{subject!non-subject} positions. Nouns\is{noun} in classes III and IV, and inanimate nouns\is{noun} in class I have retained only one (historically \isi{nominative}) morphological form for each number value (singular and plural) which is used both in the \isi{subject} and any non-subject\is{subject!non-subject} position. This is illustrated in examples (\ref{05ex10}), (\ref{05ex11}), (\ref{05ex12}) and (\ref{05ex13}) from our corpus (data represent various locations in the municipalities of Belogradchik and Trân).

\medskip

% \needspace{4\baselineskip}
Inflection class I\is{inflectional class}, \isi{masculine}. Nominative\is{nominative} and \isi{accusative} (animate nouns\is{noun} only)

\begin{exe}
	\ex \label{05ex10}
	\begin{xlist}
		\ex 
		\gll Dojde \textbf{doktor}. \\
		come[\textsc{pst.3sg}] doctor[\textsc{sg}] \\
		\trans ‘A doctor came.’ \hfill \isi{subject} NP\is{noun!phrase} \\
		\hfill Filipovtsi, Trân
		
		\ex 
		\gll Nije natovari-m pak \textbf{doktor-a}. \\
		we task-\textsc{prs.1pl} again doctor-\textsc{acc.sg} \\
		\trans ‘We task the doctor again.’ \\
		\hfill Vrabcha, Trân
		
		\ex 
		\gll Bašta mi me kara u Lom \textbf{pri} \textbf{doktor-a}. \\
		father \textsc{1sg.dat} \textsc{1sg.acc} took[\textsc{prs.3sg}] to Lom to doctor-\textsc{acc.sg} \\
		\trans ‘My father took me to the doctor in Lom.’ \\
		\hfill Chuprene, Belogradchik
	\end{xlist}
\end{exe}

% \needspace{6\baselineskip}
Inflection class II\is{inflectional class}, \isi{feminine} and \isi{masculine}. Nominative\is{nominative} and \isi{accusative}

\begin{exe}
	\ex \label{05ex11}
	\begin{xlist}
		\ex 
		\gll \textbf{Vod-a} gi e kara-l-a. \\
		water-\textsc{nom.sg} \textsc{3pl.dat} \textsc{aux.3sg} transport-\textsc{ptcp-sg.f} \\
		\trans ‘Water moved them.’ \hfill \isi{subject} NP\is{noun!phrase} \\
		\hfill Goren Chiflik, Belogradchik
		
		\ex 
		\gll Sipe-mo \textbf{vod-u}. \\
		pour-\textsc{prs.1pl} water-\textsc{acc.sg} \\
		\trans ‘We pour water.’ \hfill non-subject\is{subject!non-subject} NP\is{noun!phrase} \\
		\hfill Repljana, Belogradchik
		
		\ex 
		\gll Ide i \textbf{za} \textbf{vod-u}. \\
		go[\textsc{prs.3sg}] and for water-\textsc{acc.sg} \\
		\trans ‘And s/he is going to get water.’ \hfill non-subject\is{subject!non-subject} PP\is{prepositional!phrase} \\
		\hfill Gorni Lom, Belogradchik
	\end{xlist}
\end{exe}

Inflection class III\is{inflectional class}, \isi{feminine}. Caseless\is{case} (historically \isi{nominative})

\begin{exe}
	\ex \label{05ex12}
	\begin{xlist}
		\ex 
		\gll … \textbf{sol} se posolva. \\
		… salt[\textsc{sg}] \textsc{refl} salt[\textsc{prs.3sg}] \\
		\trans ‘… salt is being added.’ \hfill \isi{subject} NP\is{noun!phrase} \\
		\hfill Salash, Belogradchik
		
		\ex 
		\gll \textbf{Morsk-a} \textbf{sol} turi-m tri ložits-e. \\
		sea-\textsc{sg.f} salt[\textsc{sg}] pour.in-\textsc{prs.1sg} three spoon-\textsc{pl} \\
		\trans ‘I pour in three spoons of sea salt.’ \hfill non-subject\is{subject!non-subject} NP\is{noun!phrase} \\
		\hfill Ezdimirtsi, Trân
		
		\ex 
		\gll Posolva-š sâs \textbf{sol}. \\
		salt-\textsc{prs.2sg} with salt[\textsc{sg}] \\
		\trans ‘You add salt / one adds salt.’ \hfill non-subject\is{subject!non-subject} PP\is{prepositional!phrase} \\
		\hfill Salash, Belogradchik
	\end{xlist}
\end{exe}

\newpage

Inflection class IV\is{inflectional class}, \isi{neuter} and \isi{masculine}. Caseless\is{case} (historically \isi{nominative})

\begin{exe}
	\ex \label{05ex13}
	\begin{xlist}
		\ex 
		\gll … \textbf{mlek-o} se podsirva. \\
		… milk-\textsc{sg} \textsc{refl} ferment[\textsc{prs.3sg}] \\
		\trans ‘… milk is fermented.’ \hfill \isi{subject} NP\is{noun!phrase} \\
		\hfill Chuprene, Belogradchik
		
		\ex 
		\gll Ako ima-š \textbf{mlek-o}. \\
		if have-\textsc{prs.2sg} milk-\textsc{sg} \\
		\trans ‘If you have milk.’ \hfill non-subject\is{subject!non-subject} NP\is{noun!phrase} \\
		\hfill Ezdimirtsi, Trân
		
		\ex 
		\gll Dve kofišk-i \textbf{ot} \textbf{mlek-o}. \\
		two pot-\textsc{pl} of milk-\textsc{sg} \\
		\trans ‘Two pots of milk.’ \hfill non-subject\is{subject!non-subject} PP\is{prepositional!phrase} \\
		\hfill Ezdimirtsi, Trân
	\end{xlist}
\end{exe}

The clear-cut distribution of the \isi{nominative} and \isi{accusative} across \isi{subject} and non-subject\is{subject!non-subject} positions within inflection classes I and II\is{inflectional class} as presented in (\ref{05ex10}) and (\ref{05ex11}), however, does not hold in the border-region dialects\is{dialect}. These dialects\is{dialect} show a tendency to generalise the \isi{nominative} for nouns\is{noun} of inflection classes\is{inflectional class} I and II across non-subject\is{subject!non-subject} positions (making them like Standard \ili{Bulgarian} where this process has been completed in all historical inflection classes\is{inflectional class}, \citealt[146--147]{mircev1958}). This results in competing choices for case marking\is{case!marking} on non-subject\is{subject!non-subject} NPs\is{noun!phrase} and PPs\is{prepositional!phrase} as illustrated in (\ref{05ex14}) and (\ref{05ex15}) for inflection class II\is{inflectional class}:

\begin{exe}
	\ex \label{05ex14}
	\begin{xlist}
		\ex
		\gll \textbf{Vod-u} gi dava-mo. \\
		water-\textsc{acc.sg} \textsc{3pl.dat} give-\textsc{prs.1pl} \\
		\trans ‘We give them water.’ \hfill non-subject\is{subject!non-subject} NP\is{noun!phrase} (\isi{accusative}) \\
		\hfill Repljana, Belogradchik
		
		\ex 
		\gll \textbf{Vod-a} smo kara-l-i. \\
		water-\textsc{nom.sg} \textsc{aux}.\textsc{prs.1pl} carry-\textsc{ptcp-pl} \\
		\trans ‘We used to carry/haul water.’ \hfill non-subject\is{subject!non-subject} NP\is{noun!phrase} (\isi{nominative}) \\
		\hfill Ezdimirtsi, Trân
	\end{xlist}
\end{exe}

\begin{exe}
	\ex \label{05ex15}
	\begin{xlist}
		\ex 
		\gll Ne može \textbf{bez} \textbf{vod-u} da peče-š. \\
		\textsc{neg} possible without water-\textsc{acc.sg} \textsc{comp} roast-\textsc{prs.2sg} \\
		\trans ‘One can’t roast [meat] without water.’ \hfill non-subject\is{subject!non-subject} PP\is{prepositional!phrase} (\isi{accusative}) \\
		\hfill Filipovtsi, Trân
		
		\ex 
		\gll Da ne ostane \textbf{bez} \textbf{vod-a}. \\
		\textsc{comp} \textsc{neg} remain[\textsc{prs.3sg}] without water-\textsc{nom.sg} \\
		\trans ‘So that s/he won’t be left without water.’ \hfill non-subject\is{subject!non-subject} PP\is{prepositional!phrase} (\isi{nominative}) \\
		\hfill Salash, Belogradchik
	\end{xlist}
\end{exe}

From a diachronic\is{diachrony} perspective, the situation in \ili{Bulgarian} border-region dialects\is{dialect} reflects one of the final stages in the transition to uninflectedness: only two of four historical inflection classes\is{inflectional class} retain \isi{case} distinctions, classes I and II, and in only one part of the \isi{paradigm} (singular); in addition, all inanimate nouns\is{noun} of class I have already lost the ability to inflect for case\is{case!loss of case}. This is different from what we see at the other end of the dialect continuum\is{dialect!continuum}, the three groups of \ili{Serbian} dialects\is{dialect} in Section \ref{05sec3}, which illustrate the initial stages in the decomposition of the Proto-Slavic\il{Slavic!Proto-Slavic} case system\is{case!system}. This raises the question as to whether the factors that condition the loss of case distinctions\is{case!loss of case} remain stable over time, or whether they change from one stage of the process to the next. We address this question by looking into the synchronic\is{synchrony} competition between \isi{accusative} as a specific non-subject\is{subject!non-subject} \isi{case} and \isi{nominative} as a single form for all syntactic functions\is{syntactic!function}, and identify what might trigger speakers' choices where both alternatives are available. Given the limited number of examples containing animate nouns\is{noun} of class I, further analysis focuses on the nouns\is{noun} belonging to the more populous class II.

\subsection{Data}

\largerpage
The study is based on annotated transcripts of sociolinguistic interviews recorded in West Bulgaria (the dialects\is{dialect} of Belogradchik and Trân, part of the border-region dialects\is{dialect}). The fieldwork was conducted within the project ``Declining case: inflectional\is{inflection} loss in progress''\is{case!decline} by Vladimir Zhobov (University of Sofia, project consultant for \ili{Bulgarian}) and Alexander Krasovitsky in 2021 and 2022 (Figure \ref{05fig4}). Additional recordings provided by Vladimir Zhobov stem from his previous fieldwork in North-West Bulgaria.

\begin{figure}
	\centering
	\includegraphics[width=0.9\textwidth]{figures/05fig4.pdf}
	\caption{Locations of the interviews in Bulgaria. Border-region dialects\is{dialect}
		\newline
		\newline
		\textcolor{customlightblue}{\Large ●} Belogradchik sub-group \quad
		\textcolor{customyellow}{\Large ●} Trân sub-group \quad}
	\label{05fig4}
\end{figure}

Forty-six residents born between 1926 and 1950 were recorded (one interview per language consultant; the interviews range in length from forty minutes to two hours). Thirteen of the recorded consultants retain two cases\is{case}, \isi{nominative} and \isi{accusative}, in inflection class II\is{inflectional class}. These thirteen interviews make up a corpus of approximately 168,000 words containing 2292 phrases with inflection class II\is{inflectional class} nouns\is{noun!phrase} in non-subject\is{subject!non-subject} positions. All thirteen interviews demonstrate a competition of \isi{accusative} and \isi{nominative} in non-subject\is{subject!non-subject} positions, however, the probabilities vary considerably between speakers (Figure \ref{05fig5}).

\begin{figure}[H]
	\centering
	\includegraphics[width=0.8\textwidth]{figures/05fig5.pdf}
	\caption{Relative frequency of \isi{nominative} and \isi{accusative} forms in inflection class II\is{inflectional class} nouns\is{noun} in non‑subject positions (by speaker).  Corpus size: around 168, 000 tokens. 13 sociolinguistic interviews, 2292 occurrences of inflection class\is{inflectional class} II nouns\is{noun} in non-subject\is{subject!non-subject} positions. Codes below the columns identify individual speakers.}
	\label{05fig5}
\end{figure}

\is{information structure|(}
\subsection{Hypothesis: information structure conditions case selection\is{case!selection}}

Immediate observations made during the sociolinguistic interviews and further elicitation work led us to the following hypothesis. Variation\is{variation} between the \isi{nominative} and \isi{accusative} on non-subjects\is{subject!non-subject} is conditioned by information structure and pragmatic salience\is{pragmatic!salience}: nouns\is{noun} associated with new and pragmatically important information tend to lose specific non-subject\is{subject!non-subject} case marking\is{case!marking} and attract \isi{nominative}; nouns\is{noun} associated with background information and pragmatically less salient\is{pragmatic!salience} parts of an utterance tend to retain specific non-subject\is{subject!non-subject} marking and show stronger preference for the \isi{accusative}.

In what follows, we briefly consider the role of pragmatic factors\is{pragmatic!factor} in case selection\is{case!selection} within a general linguistic context (Section \ref{05sec4.3}) and then proceed to the analysis of our data (Section \ref{05sec4.4}) and a discussion (Section \ref{05sec4.5}). We argue that pragmatic conditions\is{pragmatic!condition} play a crucial role in speakers' choices between the two competing cases\is{case} which mark non-subjects\is{subject!non-subject} and in the decline of the \isi{accusative} as distinct morphological marking for non-subjects\is{subject!non-subject}. At the same time, we will show that the relationship between \isi{case} and information structure in \ili{Bulgarian} border-region dialects\is{dialect} differs significantly from that observed cross-linguistically, and offer an explanation for this.

\subsection{Information structure as a condition on case marking\is{case!marking}} \label{05sec4.3}

Information structure is claimed as a cause for variable case marking\is{case!marking} cross\hyp{}linguistically (e.g. \citealt{aikhenvald2010, iemmolo2010, valle2011, dalrymplenikolaeva2011}). First, a large body of typological research links differential object marking (DOM)\is{differential object marking} to \isi{disambiguation}: typical objects\is{object} which cannot be confused for subjects\is{subject} are less likely to receive specific grammatical marking than atypical ones, where such marking is essential for \isi{disambiguation} (\citealt{silverstein1976, comrie1977} \emph{inter alia}). With respect to information structure, this would mean that topical objects\is{object!topical} need to be marked to avoid ambiguity with subjects\is{subject}, while objects\is{object} occupying their typical slot as part of the pragmatic \isi{focus} may remain unmarked. In some languages specific non-subject\is{subject!non-subject} markers are used with highly topical objects\is{object!topical} dislocated\is{object!dislocated} from their default position (\citealt{iemmolo2010, aikhenvald2003}). Second, the degree of pragmatic prominence is another factor which may affect case marking\is{case!marking}: it has been shown for a number of DOM languages\is{differential object marking} that more prominent objects\is{object} with a higher pragmatic status are more likely to be marked for case\is{case!marking} than pragmatically less salient\is{pragmatic!salience} objects\is{object} (\citealt{aissen2003, malchukov2008}).

It is useful to assess the South Slavic situation against this broad cross-linguistic picture. The data from \ili{Bulgarian} border-region dialects\is{dialect} provide no evidence for case marking\is{case!marking} being used as a tool in contexts where atypical objects\is{object} (e.g. topicalised\is{object!topical} and dislocated\is{object!dislocated} to non-default positions or used with omitted governors\is{governor}) may be confused with subjects\is{subject} (see the analysis of the relationship between \isi{word order}, pragmatic structure\is{pragmatic!structure} and case marking\is{case!marking} in Section \ref{05sec4.4}). Thus, the situation in the border-region dialects\is{dialect} is different, in broad terms, to a situation in some of the DOM languages\is{differential object marking} (such as those considered in \citealt{aissen2003, malchukov2008}) where pragmatically more salient\is{pragmatic!salience} nouns\is{noun} attract case marking\is{case!marking} while less pragmatically salient\is{pragmatic!salience} nouns\is{noun} do not. In the \ili{Bulgarian} dialects\is{dialect} in question, case marking\is{case!marking} on non-subject\is{subject!non-subject} NPs\is{noun!phrase} is conditioned, to a significant extent, by whether they are associated with background or new information in an utterance, or with the \isi{topic}--\isi{focus} distinction. However, a significant peculiar feature of these dialects\is{dialect} is that specific non-subject\is{subject!non-subject} case marking\is{case!marking} (\isi{accusative}) is more frequent on nouns\is{noun} which are backgrounded. Nominative\is{nominative} forms (which may be used in \isi{subject} and non-subject\is{subject!non-subject} position, as examples in (\ref{05ex14}) and (\ref{05ex15}) demonstrate) are more frequent in the pragmatically salient\is{pragmatic!salience} part of an utterance. We will consider this in greater detail in the next section. It should be noted that in qualifying parts of an utterance as, \isi{topic} or \isi{focus} we follow the tradition of defining \isi{topic} as existing, background knowledge, or ``\isi{common ground}'', and \isi{focus} as part of an utterance which brings in new information and thus modifies the ``\isi{common ground}'' \citep{hoopswart2000, krifka2007}. Since originally our data collection was not intended to test the impact of information structure on case marking\is{case!marking}, and such analysis was applied later, we used broad definitions based on this opposition, and excluded from the analysis utterances where a non-ambiguous definition of information structure was not possible. Under this view, three types of utterances required special decisions: utterances with contrastive topic\is{topic!contrastive} as in (\ref{05ex16}), those with repeated focus\is{focus!repeated}, as in (\ref{05ex17}) and those with contrastive focus\is{focus!contrastive}, as in (\ref{05ex18}).


\begin{exe}
	\ex \label{05ex16}
	\textbf{Interviewer:} \\
	Kak platiš na toj kojto gledal ovci celoto leto? \\
	‘How do you pay someone who looked after sheep the whole summer?’
	
	\medskip
	
	\textbf{Consultant:} \\
	\gll Ne ima-l-o takâv. \textbf{U} \textbf{naš-a-ta} \textbf{bačij-a} ne ima-l-o ovčar-e. \\
	\textsc{neg} \textsc{exist-ptcp-sg.n} such at \textsc{1pl}-\textsc{nom.sg-def.nom.sg} village.sheep-\textsc{nom.sg} \textsc{neg} \textsc{exist-ptcp-sg.n} shepherd-\textsc{pl} \\
	\trans ‘There was no such person. For our village sheep, there were no shepherds.’ \hfill (contrastive topic\is{topic!contrastive})
\end{exe}


\begin{exe}
	\ex \label{05ex17}
	\gll A bab-a mi ode-še… pase-še \textbf{krav-a}… Ja njakolko pât otiš-l-a sâs nju*, ama poveče ona si e pas-l-a \textbf{krav-u}. \\
	and grandmother-\textsc{nom.sg} \textsc{1sg.dat} go-\textsc{impf.3sg} graze-\textsc{impf.3sg} cow-\textsc{nom.sg} \textsc{1sg.nom} several time go-\textsc{ptcp-sg.f} with \textsc{3sg.acc.f} but more \textsc{3sg.nom.f} \textsc{part} \textsc{aux.prs.3sg} graze\textsc{-ptcp-sg.f} cow\textsc{-acc.sg} \\
	\trans ‘And my grandmother used to go… grazed the cow. I went with her a few times but mostly she grazed the cow.’ \hfill (repeated focus\is{focus!repeated})
\end{exe}

{\footnotesize *Dialectal\is{dialect} form used in the border-region and some other dialects\is{dialect} corresponding with the Standard \ili{Bulgarian} \textit{neja} ‘she [\textsc{3sg.acc.f}]’ (cf. \citealt[82]{todorov2002}).}

\begin{exe}
	\ex \label{05ex18}
	\textbf{Interviewer:} \\
	A kâde go slagate kopâra? \\
	‘Where (i.e. in which dishes) do you add dill?’
	
	\medskip
	
	\textbf{Consultant:} \\
	\gll U tarator tura-m mnogo. Daže i \textbf{u} \textbf{salat-a} si tura-m. \\
	in tarator[\textsc{sg}] put-\textsc{1sg} much even and in salad-\textsc{nom.sg} \textsc{part} put-\textsc{1sg} \\
	\trans ‘I put a lot of dill in tarator. I even put it in salad.’ \hfill (contrastive focus\is{focus!contrastive})
\end{exe}

Following the approach taken in this study, we grouped together utterances where non-subject\is{subject!non-subject} nouns\is{noun} appear in the repeated focus\is{focus!repeated}, as in (\ref{05ex17}), and in the \isi{topic} on the one hand, and those where non-subject\is{subject!non-subject} nouns\is{noun} are part of the contrastive topic\is{topic!contrastive}, as in (\ref{05ex16}), and of the \isi{focus}. The rationale behind both decisions is based on information novelty. In the former case, information in the repeated focus\is{focus!repeated} is already brought into the \isi{discourse} and cannot be defined as changing the `\isi{common ground}'. In the latter case, contrastive topic\is{topic!contrastive} turns out to be very close to \isi{focus}, in that the contrast itself contributes to existing knowledge and thus modifies the `\isi{common ground}'. The similarity of \isi{focus} and contrastive topic\is{topic!contrastive} has been demonstrated in a number of studies within Alternative Semantics, where both concepts are understood as parts of an utterance which provide new information by singling out existing alternatives (\citealt[75--77]{rooth1992}; \citealt[64--68]{buring2016}). A small number of examples were annotated as contrastive focus\is{focus!contrastive}, as in (\ref{05ex18}). Contrastive focus\is{focus!contrastive} is not different from \isi{focus} in that it, like \isi{focus}, provides new information and modifies the `\isi{common ground}', under some specific conditions such as polarity of unexpectedness (e. g. \citealt[348]{zimmermann2008}; \citealt[117--126]{goodhue2022}). For this reason, in the analysis presented in Section \ref{05sec4.4} we did not distinguish between the two types of \isi{focus}. However, a possible effect of contrastive focus\is{focus!contrastive} will be discussed separately in the account presented in Section \ref{05sec4.5}.

We will consider the effect of information structure on non-subject\is{subject!non-subject} case marking\is{case!marking} with respect to two other factors: \isi{word order} and phrase\is{noun!phrase} type (NP\is{noun!phrase} vs. PP\is{prepositional!phrase}).

Previous research has suggested that the diachronic\is{diachrony} loss of case marking\is{case!marking} is accompanied by compensatory freezing of \isi{word order} as a means of distinguishing grammatical roles\is{grammatical role} \citep{blake2004}. Among other things, this implies that if an \isi{object} \isi{noun} occurs outside of its expected position, it will be more likely to preserve its original case marking\is{case!marking}. This straightforward complementarity between fixed \isi{word order} and case marking\is{case!marking} does not necessarily hold (cf. \citealt{allen2006} for Old English\il{English!Old} and \citealt{detges2009} for Old French\il{French!Old}); in fact, the opposite correlation may occur. Data from \ili{Japanese} show that \isi{word order} alternations may loosen grammatical relationships within a sentence so that dislocated objects\is{object!dislocated} may be \textit{less} likely to receive expected case marking\is{case!marking} than their counterparts in neutral positions within a basic constituent structure (cf. \ili{Japanese} scrambling of \isi{accusative}-marked objects\is{object} to a position where only the \isi{nominative} may be assigned, \citealt{kasai2018}).

The effect of phrase type (noun phrases\is{noun!phrase} vs. prepositional phrases\is{prepositional!phrase}) on case selection\is{case!selection} and decline of one of the competing cases\is{case} was surveyed in Section \ref{05sec3}. Here we will only reiterate the observation made above that the correlation between the use of prepositions\is{preposition} and the tendency to maintain an original \isi{case} form has been shown to be strong for the \isi{dialect} types under study. Therefore, we found it important to disentangle different conditions on case marking\is{case!marking} in order to assess their effect individually without interference from other conditions. To do so, we will analyse the effect of \isi{topic}-\isi{focus} distinction separately for different types of \isi{word order} and different phrase types (NP\is{noun!phrase} and PP\is{prepositional!phrase})\is{information structure|)}.

\subsection{The effect of \isi{topic} and \isi{focus} on case marking\is{case!marking} in \ili{Bulgarian} border-region dialects\is{dialect}} \label{05sec4.4}

To investigate the effect of \isi{information structure} on competing choices in case marking\is{case!marking}, a sub-corpus of approximately 90,000 words (six sociolinguistic interviews) was annotated with respect to \isi{information structure}. The percentage of \isi{nominative} on non-subjects\is{subject!non-subject} in the six included interviews is as follows: 7\% (BEL1-12), 20\% (BEL1-15), 23\% (BEL1-31), 42\% (KS1-09), 55\% (KS1-08), 57\% (KS1-18). A dataset of 956 utterances containing inflection class II\is{inflectional class} nouns\is{noun} in non-subject\is{subject!non-subject} positions and allowing unambiguous assignment of the five values (\isi{topic}, contrastive topic\is{topic!contrastive}, \isi{focus}, contrastive focus\is{focus!contrastive}, repeated focus\is{focus!repeated}) was extracted. The data were combined for (i) \isi{topic} and repeated focus\is{focus!repeated}, and for (ii) \isi{focus}, contrastive focus\is{focus!contrastive} and contrastive topic\is{topic!contrastive}, as discussed in Section \ref{05sec4.3}. The analysis presented in this section is based on these two subsets.

\begin{figure}
	\centering
	\includegraphics[width=0.52\textwidth]{figures/05fig6.pdf}
	\caption{Nominative\is{nominative} and \isi{accusative} on non-subject\is{subject!non-subject} nouns\is{noun}}
	\label{05fig6}
\end{figure}

As Figure \ref{05fig6} demonstrates, there is a significant difference in the frequency of the \isi{nominative} and \isi{accusative} across the two subsets. The \isi{accusative} dominates in \isi{topic} (\ref{05ex19a}) or repeated focus\is{focus!repeated} position (\ref{05ex19b}): 78\% \isi{accusative} forms. Frequency of the \isi{accusative} falls significantly in \isi{focus} (\ref{05ex19c}) and contrastive topic\is{topic!contrastive} positions (\ref{05ex19d}): 51\%. In other words, non-subject\is{subject!non-subject} nouns\is{noun} are more likely to take the \isi{accusative} in pragmatically less salient\is{pragmatic!salience} positions; pragmatically more salient\is{pragmatic!salience} parts of a sentence show a stronger preference for the \isi{nominative} on non-subjects\is{subject!non-subject}.

\begin{exe}
% 	\needspace{6\baselineskip}
	\ex
	\begin{xlist}
		\ex \label{05ex19a}
		\gll I slam-a-ta otzad. I otvârlj-u \textbf{slam-u-tu} nastranu. \\
		and straw-\textsc{nom.sg-def.nom.sg} behind and toss-\textsc{prs.3pl} straw-\textsc{acc.sg-def.acc.sg} aside \\
		\trans ‘And the straw is behind. And they toss the straw aside.’ \hfill (\isi{topic})
		
		\ex \label{05ex19b}
		\gll Čorapete ot \textbf{dvojk-a} si plete-mo. \\
		socks from double.thread-\textsc{nom.sg} \textsc{part} weave-\textsc{prs.1pl} \\
		\trans ‘We weave socks with double threads.’ \hfill (\isi{focus}) \\[1ex]
		\gll I ot \textbf{dvojk-u} i ot trojk-u sam ple-l-a. \\
		and from double.thread-\textsc{acc.sg} and from triple.thread-\textsc{acc.sg} be.\textsc{prs.1sg} weave-\textsc{ptcp-sg.f} \\
		\trans ‘I have weaved both with double and triple threads.’ \hfill (repeated focus\is{focus!repeated})
		
		\ex \label{05ex19c}
		\gll Ako nema rek-a nosi-mo \textbf{vod-a}. \\
		if \textsc{exist.neg} river-\textsc{nom.sg} carry-\textsc{prs.1pl} water-\textsc{nom.sg} \\
		\trans ‘If there is no river, we carry water.’ \hfill (\isi{focus})
		
		\ex \label{05ex19d}
		\textbf{Interviewer:} \\
		S râka li pravete? \\
		‘Do you do it manually (``with a hand")?’ \\
		
		\medskip
		
% 		\needspace{3\baselineskip}
		\textbf{Consultant:} \\
		\gll S râk-u, a s kvo? \\
		with hand-\textsc{acc.sg} and with what \\
		\gll \textbf{S} \textbf{lev-a-ta} tegli-š otgore, \\
		with left-\textsc{nom.f.sg-def.nom.sg} pull-\textsc{prs.2sg} from.above \\
		\gll \textbf{s} \textbf{desn-a-ta} vârti-š vreten-o-to. \\
		with right-\textsc{nom.f.sg-def.nom.sg} rotate-\textsc{prs.2sg} spindle-\textsc{sg-def.sg.n} \\
		\trans ‘Manually, how else? With your left hand you pull it from above, with your right hand you rotate the spindle.’ \hfill (contrastive topic\is{topic!contrastive})
	\end{xlist}
\end{exe}

We assessed the effect of \isi{information structure} on case marking\is{case!marking} in connection with two potentially relevant factors, phrase type (presence or absence of a \isi{preposition}) and \isi{word order}. As shown above, phrase type proved to have an effect on case marking\is{case!marking} in the \ili{Serbian} dialects\is{dialect} at an earlier stage of case decline\is{case!decline} (Section \ref{05sec3}). Word order\is{word order}, although not independent of \isi{information structure}, may still vary within the same \isi{topic}--\isi{focus} structure, therefore this factor was also disentangled from the other two factors under study. In \ili{Bulgarian}, direct\is{object!direct} and indirect objects\is{object!indirect} under basic default \isi{word order} occur postverbally\is{verb}. Non-default ordering\is{word order} (e.g. preverbal\is{verb} objects\is{object}) is motivated by a variety of conditions, which may or may not be related to \isi{information structure} \citep{georgieva1974, georgieva1983, dyer1992}. Thus, topicalised objects\is{object!topical} are normally fronted, but this is not a universal rule, and under appropriate contextual conditions topicalised objects\is{object!topical} may occur postverbally\is{verb} (\ref{05ex20a}).
	
\begin{exe}
	\ex
	\begin{xlist}
	\ex \label{05ex20a}
	\textbf{Consultant 1:} \\
	\gll Kâde utiva-š? \\
	Where go-\textsc{prs.2sg}? \\
	
	\medskip
	
	\textbf{Consultant 2:} \\
	\gll Po rabot-a, majk-o. \\
	for work-\textsc{nom.sg} mother-\textsc{voc} \\
	
	\medskip
	
	\textbf{Consultant 1:} \\
	\gll Ah, mrazj-a az \textbf{tazi} \textbf{rabot-a}! \\
	ah hate-\textsc{prs.1sg} \textsc{1sg.nom} this.\textsc{nom.sg.f} work-\textsc{nom.sg} \\
	\trans ‘C1: Where are you going? C2: On business, mother. C1: Oh I hate this business!’
	
\medskip

\item[]  % this is a dummy (blank) item to be able to insert a paragraph inside an example environment (in an attempt to have 20a, 20b numbering, split with a paragraph)
\noindent\hspace*{-4.5em}\parbox{\textwidth}{On the other hand, as expected in SVO languages, objects\is{object} in \isi{focus} normally occur postverbally\is{verb}, but strong logical stress may result in emphatic \isi{word order}, under which an \isi{object} occurs in preverbal\is{verb} position (\ref{05ex20b}).}

\medskip

	\ex \label{05ex20b}
	\gll \textbf{Cjal-a-ta} \textbf{planin-a} vdišva-m. \\
	whole-\textsc{nom.sg.f-def.nom.sg.f} mountain-\textsc{nom.sg} inhale-\textsc{prs.1sg} \\
	\trans ‘I inhale all this mountain range.’ (adapted from \citealt[284]{georgieva1983}).
	\end{xlist}
\end{exe}

In addition to postverbal\is{verb} and preverbal\is{verb} ordering\is{word order}, we found a significant number of elliptic constructions\is{construction}, for example, objects\is{object} with a missing governing \isi{verb}, or interruptions where a \isi{governor} and a governee stand far apart from each other in different syntagms. To see whether the absence of an overt \isi{governor} has an effect on case marking\is{case!marking}, such phrases (tagged as ``isolated'') were considered separately.

Therefore, in order to disentangle the contribution of potentially conditioning factors, we assessed the effect of \isi{information structure} separately for six different combinations of \isi{word order} and phrase type:

\begin{exe}
	\ex
	\begin{xlist}
		\ex \label{05ex21a} NP\is{noun!phrase}, postverbal\is{verb} \\
		\gll Ne smo slaga-l-i \textbf{vod-u}. \\
		\textsc{neg} \textsc{aux.prs.1pl} put-\textsc{ptcp-pl} water-\textsc{acc.sg} \\
		\trans ‘We added no water.’
		
		\ex \label{05ex21b} NP\is{noun!phrase}, preverbal\is{verb} \\
		\gll \textbf{Krav-a} smo ima-l-i. \\
		cow-\textsc{nom.sg} \textsc{aux.prs.1pl} have-\textsc{ptcp-pl} \\
		\trans ‘We had a cow.’
		
		\ex \label{05ex21c} NP\is{noun!phrase}, isolated \\
		\gll Ne sâm ima-l-a bab-u, \\
		\textsc{neg} \textsc{aux.1sg} have-\textsc{ptcp-sg.f} grandmother-\textsc{acc.sg} \\
		\gll samo \textbf{majk-a}. \\
		only mother-\textsc{nom.sg} \\
		\trans ‘I didn’t have a grandmother, just a mother.’
		
		\ex \label{05ex21d} PP\is{prepositional!phrase}, postverbal\is{verb} \\
		\gll I ojde-mo večerom \textbf{na} \textbf{večerj-u}. \\
		and go-\textsc{prs.1pl} in.the.evening to dinner-\textsc{acc.sg} \\
		\trans ‘And in the evening we will go to have dinner.’
		
		\largerpage[-1.5]

% 		\needspace{4\baselineskip}
		\ex \label{05ex21e} PP\is{prepositional!phrase}, preverbal\is{verb} \\
		\gll \textbf{Na} \textbf{bašt-a} \emph{mi} pomaga-x. \\
		to father-\textsc{nom.sg} \textsc{1sg.dat} help-\textsc{aor.1sg} \\
		\trans ‘I helped my father.’
		
% 		\needspace{3\baselineskip}
		\ex \label{05ex21f} PP\is{prepositional!phrase}, isolated \\
		\textbf{Interviewer:} 
		\\ Ima li češma? \\
		\trans ‘Is there a tap?’ \\
		\textbf{Consultant:} \\
		\gll Otzad češm-a-ta. \textbf{Na} \textbf{kuxničk-u}. \\
		in.the.rear tap-\textsc{nom.sg-def.nom.sg} in kitchen-\textsc{acc.sg} \\
		\trans ‘The tap is in the rear. In the kitchen.’
	\end{xlist}
\end{exe}

Relative frequencies of the two alternative cases\is{case}, \isi{nominative} and \isi{accusative}, on non-subject\is{subject!non-subject} nouns\is{noun} with respect to the six combinations of conditions are presented in Figure \ref{05fig7} (for NPs\is{noun!phrase}) and Figure \ref{05fig8} (for PPs\is{prepositional!phrase}).

\begin{figure}[ht]
	\centering
	\begin{subfigure}{0.462\textwidth}
		\includegraphics[width=\linewidth]{figures/05fig7a.pdf}
		\caption{POSTVERBAL NP\is{noun!phrase}}
	\end{subfigure}
	\hfill
	\begin{subfigure}{0.498\textwidth}
		\includegraphics[width=\linewidth]{figures/05fig7b.pdf}
		\caption{PREVERBAL NP\is{noun!phrase}}
	\end{subfigure}

	\vspace{1em}
	
	\begin{subfigure}{0.5\textwidth}
		\includegraphics[width=\linewidth]{figures/05fig7c.pdf}
		\caption{ISOLATED NP\is{noun!phrase}}
	\end{subfigure}

	\vspace{1em}

	% Legend
	\begin{tabular}{llll}
		\textcolor{customdarkblue}{\rule{0.5em}{0.5em}} & \textsc{nom} & \textcolor{customorange}{\rule{0.5em}{0.5em}} & \textsc{acc} \\
	\end{tabular}

	\caption{Preverbal\is{verb}, postverbal and isolated NPs\is{noun!phrase} in \isi{topic} and \isi{focus} positions}
	\label{05fig7}
\end{figure}




\begin{figure}[ht]
	\centering
	\begin{subfigure}{0.462\textwidth}
		\includegraphics[width=\linewidth]{figures/05fig8a.pdf}
		\caption{POSTVERBAL PP\is{prepositional!phrase}}
	\end{subfigure}
	\hfill
	\begin{subfigure}{0.498\textwidth}
		\includegraphics[width=\linewidth]{figures/05fig8b.pdf}
		\caption{PREVERBAL PP\is{prepositional!phrase}}
	\end{subfigure}
	
	\vspace{1em}
	
	\begin{subfigure}{0.51\textwidth}
		\includegraphics[width=\linewidth]{figures/05fig8c.pdf}
		\caption{ISOLATED PP\is{prepositional!phrase}}
	\end{subfigure}

	\vspace{1em}
	
	% Legend
	\begin{tabular}{llll}
		\textcolor{customdarkblue}{\rule{0.5em}{0.5em}} & \textsc{nom} & \textcolor{customorange}{\rule{0.5em}{0.5em}} & \textsc{acc} \\
	\end{tabular}

	\caption{Preverbal\is{verb}, postverbal and isolated NPs\is{noun!phrase} in \isi{topic} and \isi{focus} positions}
	\label{05fig8}
\end{figure}

A key outcome of this analysis is the uniform pattern which we found across all six data subsets (factor combinations): while pragmatically more important parts of an utterance (\isi{focus} or contrastive topic\is{topic!contrastive}) demonstrate a strong tendency to generalize the \isi{nominative} on non-subject\is{subject!non-subject} nouns\is{noun}, the \isi{accusative} as a distinct non-subject\is{subject!non-subject} \isi{case} is better preserved in pragmatically less salient\is{pragmatic!salience} positions (\isi{topic} and repeated focus\is{focus!repeated}). In the light of the findings reported above for \ili{Serbian} dialects\is{dialect}, it is particularly striking that structural conditions, that is, the absence or presence of a \isi{preposition}, seem to have no impact on speakers' choices, and neither enhance the use of a distinct non-subject\is{subject!non-subject} \isi{case} form (\isi{accusative}), nor contribute to its loss and the generalisation of an undifferentiated morphological form (\isi{nominative}). In other words, non-subject\is{subject!non-subject} NPs\is{noun!phrase} and PPs\is{prepositional!phrase} in terms of their case marking\is{case!marking} show similar sensitivity to \isi{information structure}. With respect to \isi{word order}, a strong effect of \isi{information structure} is found for all three linear structures: postverbal\is{verb}, preverbal\is{verb} and isolated positions. In all three data subsets, the \isi{accusative} is better retained in \isi{topic} and repeated focus\is{focus!repeated} position, and is less frequent in \isi{focus} and contrastive topic\is{topic!contrastive} position where it is replaced by the \isi{nominative}.

The statistical significance of the results presented in Figure \ref{05fig7} and Figure \ref{05fig8} was assessed by ANOVA analysis. The effect of each of the three factors, phrase type, \isi{information structure} and \isi{word order}, was considered as a sole factor, and in connection with the other two factors (the null hypothesis assumed that there was no effect of either of these factors on case selection\is{case!selection}). The outcome of this analysis is presented in Table \ref{05tab4}.

\begin{table}[ht]
	\centering
	\caption{The effect of phrase type, \isi{word order} and \isi{information structure} on case marking\is{case!marking} in \ili{Bulgarian} border-region dialects\is{dialect}}
	\label{05tab4}
	\begin{tabularx}{\textwidth}{XXX}
		\lsptoprule
		Phrase type (NP\is{noun!phrase} vs. PP\is{prepositional!phrase}) & Word order\is{word order} & Information structure\is{information structure} \\
		\midrule
		NO EFFECT & WEAK EFFECT (?) & STRONG EFFECT \\
		& & \\
		p = 0.35 in connection with WO and IS & p = 0.18 in connection with WO and phrase type & p < 0.001 under any condition(s) or taken separately \\
		p = 0.29 taken separately & p < 0.001 taken separately & \\
		\lspbottomrule
	\end{tabularx}
\end{table}

The results of the statistical analysis presented in Table \ref{05tab4} confirm the observation that phrase type (NP\is{noun!phrase} or PP\is{prepositional!phrase}), whether taken separately or in connection with \isi{word order} and \isi{information structure}, has no effect on case marking\is{case!marking} on non-subjects\is{subject!non-subject} (NP: p=0.29 and PP: p=0.35, respectively). On the other hand, we see a strong effect of \isi{information structure}, either considered on its own (as in Figure \ref{05fig6}), or in connection with phrase type and \isi{word order}, for the sub-sets presented in Figures \ref{05fig7} and \ref{05fig8} (p \textless{} 0.001 either taken separately or in connection with other conditions). We cannot, however, arrive at such an unequivocal conclusion on the role of \isi{word order}. Within each of the six sub-sets presented in Figures \ref{05fig7} and \ref{05fig8}, \isi{word order} does not reach clear statistical significance (p=0.18). Taken as a sole factor without considering other factors (structure of the phrase and its affiliation with different pragmatic roles) postverbal\is{verb} positions seem to enhance the use of the \isi{accusative}, and preverbal\is{verb} positions favour the \isi{nominative} (p \textless{} 0.001). However, it should be noted that the latter outcome may reflect a natural bias in the data. Non-subject\is{subject!non-subject} postverbal\is{verb} \isi{focus} NPs\is{noun!phrase} are more frequent than preverbal\is{verb} ones. Given that \isi{focus} tends to attract \isi{nominative} marking according to our data, this may significantly bias the analysis of the role of \isi{word order} as a sole factor in case marking\is{case!marking}. Therefore, we think it important that \isi{word order}, when taken as a factor in combination with \isi{information structure}, fails to reach unambiguous statistically significant values. Firm conclusions, however, could only be reached from data normalised with respect to different types of \isi{word order} and different pragmatic roles, which would require a larger corpus than we currently have.

\is{case!decline|(}
\subsection{Discussion: pragmatic factors\is{pragmatic!factor} in case decline} \label{05sec4.5}

The reported findings leave little doubt that \isi{information structure} has a crucial effect on \isi{variation} in case marking\is{case!marking} in \ili{Bulgarian} border-region dialects\is{dialect}. What might this synchronic\is{synchrony} \isi{variation} tell us about the \isi{diachrony} of case decline? We hypothesise that the historical competition of case marking\is{case!marking} here is motivated by pragmatic salience\is{pragmatic!salience}. What we observe, however, is not a tight coupling between a pragmatic role and grammatical form, but rather a tendency for caseless\is{case} forms (historically the \isi{nominative}) to be generalised first across the pragmatically most salient\is{pragmatic!salience} positions, and only then to spread further to other parts of the utterance. Key evidence supporting this view comes from the comparison of different parts of the subcorpus that have been annotated for \isi{information structure}. Preferences for case marking\is{case!marking} on non-subject\is{subject!non-subject} nouns\is{noun} in this subcorpus vary enormously across speakers, ranging from a total of 7\% to 57\% of \isi{nominative} forms. We divided the subcorpus into two parts depending on the by-speaker frequency of \isi{nominative} forms. The first part includes the three most conservative speakers (from 7\% to 23\% of \isi{nominative} forms in non-subject\is{subject!non-subject} positions). The second part includes interviews recorded from more innovative speakers with an overall frequency of \isi{nominative} forms between 42\% and 57\% (T-test confirms the significance of the difference between the two groups with respect to case variation\is{case!variation}: p\textless0.0001, t=8.4).

The analysis presented below shows which parts of the \isi{information structure} are most susceptible to the loss of \isi{case} distinctions at a very early stage of this process. We looked at the data from the three most conservative speakers in order to identify conditions which make these speakers deviate from their dominant pattern, \isi{accusative} marking of non-subjects\is{subject!non-subject}. We assessed the impact of each pragmatic condition\is{pragmatic!condition} for which our data are annotated. Unlike the findings from the whole subcorpus analysed in the previous section (six interviews), our results here show that in the three most conservative individual dialects\is{dialect}, the \isi{accusative} is strongly associated not only with the background or repeated information (\isi{topic} and repeated focus\is{focus!repeated}), but also with new information (\isi{focus}). If, however, we look separately at the pragmatically most salient\is{pragmatic!salience} NPs\is{noun!phrase} and PPs\is{prepositional!phrase}, contrastive topic\is{topic!contrastive} and contrastive focus\is{focus!contrastive}, we find that these pragmatic conditions\is{pragmatic!condition} favour the \isi{nominative} (Table \ref{05tab5}, illustrated by examples in (\ref{05ex22})).

\largerpage[1.5]

\begin{table}[ht]
	\centering
	\caption{Distribution of \isi{nominative} and \isi{accusative} forms with non-subject\is{subject!non-subject} nouns\is{noun} for the three most conservative individual dialects\is{dialect}. Raw numbers and \% of \isi{nominative}}
	\label{05tab5}
	\begin{tabularx}{0.87\textwidth}{lrrr}
		\lsptoprule
		& \textsc{\isi{accusative}} & \textsc{nominative} & \% \textsc{nominative} \\
		\midrule
		Topic\is{topic} & 17 & 0 & 0\% \\
		Repeated focus\is{focus!repeated} & 13 & 4 & 24\% \\
		Focus\is{focus} & 104 & 18 & 15\% \\
		Contrastive topic\is{topic!contrastive} & 0 & 3 & 100\% \\
		Contrastive focus\is{focus!contrastive} & 0 & 2 & 100\% \\
		\midrule
		Total & 134 & 27 & 17\% \\
		\lspbottomrule
	\end{tabularx}
\end{table}

\begin{exe}
	\ex \label{05ex22}
	\begin{xlist}
		\ex
		\textbf{Interviewer:} \\
		Kak platiš na toj kojto gledal ovci celoto leto? \\
		‘How do you pay someone who looked after sheep the whole summer?’
		
		\medskip
		
		\textbf{Consultant:} \\
		\gll Ne ima-l-o takâv. \textbf{U} \textbf{naš-a-ta} \textbf{bačij-a} ne ima-l-o ovčar-e. \\
		\textsc{neg} \textsc{exist-ptcp-sg.n} such at \textsc{1pl}-\textsc{nom.sg.f-def.nom.sg.f} village.sheep-\textsc{nom.sg} \textsc{neg} \textsc{exist-ptcp-sg.n} shepherd-\textsc{pl} \\
		\trans ‘There was no such person. For our village sheep, there were no shepherds.’ \hfill (contrastive topic\is{topic!contrastive})
		
		\ex
		\textbf{Interviewer:} \\
		Da mu platiš posle? \\
		‘Do you pay him afterwards?’
		
		\medskip
		
		\textbf{Consultant:} \\
		Ne, da ne mu platiš. \\
		‘No, you do not pay him.’
		
		\gll Odi-l-i smo \textbf{na} \textbf{zared-a}. \\
		go-\textsc{ptcp-pl} \textsc{aux.prs.1pl} on work.for.smb-\textsc{nom.sg} \\
		\trans ‘We went to work for him (in exchange for his favour).’ \\
		\hfill (contrastive focus\is{focus!contrastive})


\medskip

\item[]  % this is a dummy (blank) item to be able to insert a paragraph inside an example environment (in an attempt to have 20a, 20b numbering, split with a paragraph)
\noindent\hspace*{-4.5em}\parbox{\textwidth}{Under non-contrastive focus\is{focus!contrastive} we can expect the \isi{accusative} (\ref{05ex22c}):}

\medskip

	\ex \label{05ex22c}
	\gll  I posle ide-mo \textbf{za} \textbf{nevest-u-tu}. \\
	and after go-\textsc{prs.1pl} for bride\textsc{-acc.sg-def.acc.sg}  \\
	\trans ‘And then we go to pick up the bride.’ \\
	\hfill (non-contrastive focus\is{focus!contrastive}, listing consecutive actions)
	\end{xlist}
\end{exe}

The data from the three most innovative dialects\is{dialect} reveal the spread of \isi{nominative} from the most salient pragmatic conditions\is{pragmatic!condition}, such as contrastive topic\is{topic!contrastive} and contrastive focus\is{focus!contrastive}, to the parts of the utterance containing new information in general (as demonstrated in Table \ref{05tab6}).

\begin{table}[ht]
	\centering
	\caption{Distribution of \isi{nominative} and \isi{accusative} forms with non-subject\is{subject!non-subject} nouns\is{noun} for the three most innovative individual dialects\is{dialect}. Raw numbers and \% of nominative}
	\label{05tab6}
	\begin{tabularx}{0.85\textwidth}{lrrr}
		\lsptoprule
		& \textsc{\isi{accusative}} & \textsc{nominative} & \% \textsc{nominative} \\
		\midrule
		Topic\is{topic} & 70 & 33 & 32\% \\
		Repeated focus\is{focus!repeated} & 69 & 10 & 13\% \\
		Focus\is{focus} & 272 & 326 & 55\% \\
		Contrastive topic\is{topic!contrastive} & 1 & 4 & 80\% \\
		Contrastive focus\is{focus!contrastive} & 1 & 9 & 90\% \\
		\midrule
		Total & 413 & 382 & 48\% \\
		\lspbottomrule
	\end{tabularx}
\end{table}

The data for a subset of the most conservative speakers in the corpus, however small, may be taken as an indication of a diachronic\is{diachrony} tendency when compared to the whole corpus. While \isi{accusative} dominates in the speech of the three conservative speakers overall, the pragmatically most salient\is{pragmatic!salience} parts of an utterance favour the \isi{nominative} on non-subject\is{subject!non-subject} nouns\is{noun}. This process has further developed with more innovative speakers with whom \isi{nominative} has spread further into the \isi{focus} position, where it dominates, and into \isi{topic} position, though here it is still in the minority.

The fact that pragmatics play a crucial role in speakers' choices in \ili{Bulgarian} border-region dialects\is{dialect} may be due to both language-internal and external factors. As we noted earlier in Section \ref{05sec4}, nouns\is{noun} of inflection classes I and II\is{inflectional class} are the only domain where \isi{case} distinctions are still preserved, while nouns\is{noun} of the other two classes (and adjectival forms\is{adjectival construction} which agree with them) are uninflected. This contributes to the marginalisation of \isi{case} and restricts \isi{accusative} to less salient positions. This process is supported by external factors, namely exposure to language varieties without \isi{case} distinctions on a daily basis (younger speakers, TV, written language).This can marginalise \isi{accusative} forms, restricting them to less highlighted parts of an utterance\is{case!decline|)}.

\section{Conclusion}

The ongoing loss of case\is{case!loss of case} in \ili{Serbian} and \ili{Bulgarian} dialects\is{dialect} results in the competition between two forms, one of them more specific and one of them more general. How exactly this competition is manifested depends on the particular case system\is{case!system} it occurs in. In more conservative varieties that retain in principle a six-\isi{case} system (\ili{Serbian}), competition is between a more specific \isi{case} (e.g. the \isi{genitive}) and a general oblique\is{oblique!general} (\isi{accusative}). In more innovative varieties (\ili{Bulgarian}), where case loss\is{case!loss of case} is more advanced, the competition is between the general oblique\is{oblique!general} form (\isi{accusative}) and an uninflected stem. A particularly interesting result of our study has been to uncover key factors underlying these systems of competing case marking\is{case!marking}. Broadly construed, they match what has been observed cross-linguistically, but the principles that govern them are different. Elsewhere, prior work has shown that in case systems\is{case!system} which allow for more explicit and less explicit marking for a given syntactic role\is{syntactic!role}, more explicit case marking\is{case!marking} has two major functions: (i) indicating the syntactic role\is{syntactic!role} of nominals\is{nominal} in contexts where this is not otherwise obvious, and (ii) highlighting nominals\is{nominal} that are semantically or pragmatically prominent. In the \ili{Serbian} and \ili{Bulgarian} dialects\is{dialect} investigated here, the factors conditioning the competition of cases\is{case} appear to be the exact opposite.

First, contexts where the syntactic role\is{syntactic!role} of a \isi{nominal} is unambiguously signalled, as in prepositional phrases\is{prepositional!phrase}, in fact provide a more favourable environment for the preservation of the original \isi{case} forms. We have illustrated this with the data from \ili{Serbian} dialects\is{dialect} where the original \isi{genitive} \isi{case} forms alternate with the innovative \isi{accusative} forms. We analysed three \isi{construction} types: prepositional constructions\is{prepositional!construction}, adnominal constructions\is{adnominal construction}, and constructions\is{construction} with a \isi{quantifier}. In all the dialects\is{dialect} of the sample, prepositional constructions\is{prepositional!construction} preserved the original \isi{case} forms better than the other two. The three constructions\is{construction} vary with respect to the strength of connection between their elements. While the use of a \isi{noun} (such as `cup' or `top') or a \isi{quantifier} does not necessarily entail the subsequent use of a \isi{noun}, a \isi{preposition} obligatorily governs a noun phrase\is{noun!phrase}, marked with the \isi{genitive} in our case. Thus, the established association of a form with its apparent context turns out to be a key factor which determines speakers' choices and contributes to the preservation of the original \isi{case} forms.

Second, higher pragmatic salience\is{pragmatic!salience}, rather than triggering \isi{case} distinctions, as has been widely observed cross-linguistically, in fact leads to the opposite result in the \ili{Bulgarian} border-region dialects\is{dialect}. These dialects\is{dialect} retain at most two cases\is{case}, \isi{nominative} and \isi{accusative}. The \isi{accusative} is normally used for non-subjects\is{subject!non-subject}, but this is not always so: they may also take the \isi{nominative}, making them morphologically indistinguishable from subjects\is{subject} and thus effectively uninflected. This \isi{variation} is conditioned by \isi{information structure}: a \isi{subject} (\isi{nominative}) vs. non-subject\is{subject!non-subject} (\isi{accusative}) distinction is better maintained where background (given) information is being presented, while pragmatically salient\is{pragmatic!salience} parts of an utterance show a strong tendency to eliminate this distinction. Additional factors that play a role in limiting the use of the \isi{accusative} are the fact that only two classes of nouns\is{noun} are affected (inflection classes I and II\is{inflectional class}), and the overall marginalisation of \isi{case} distinctions on nouns\is{noun} in the context of the prestige standard language, which lacks \isi{nominal} \isi{case}.

The fine-grained analysis of the competing patterns of case marking\is{case!marking} provided here contributes to our understanding of factors that condition the loss of inflectional\is{inflection} case\is{case!loss of case}. In particular, different factors appear to be activated at different stages of the change. A synchronic\is{synchrony} cross-\isi{dialect} analysis therefore allows us to uncover the intricate detail of this historical process\il{Slavic!South|)}.

 
 
\section*{Abbreviations}
For interlinear glosses, we follow the conventions of the Leipzig Glossing Rules (\url{https://www.eva.mpg.de/lingua/pdf/Glossing-Rules.pdf}) with the following additions:

\medskip

	\begin{tabularx}{.45\textwidth}{lQ}
	\textsc{aor} & aorist \\
	\textsc{exist} & existential \\
	\textsc{impf} & imperfect \\
	\end{tabularx}



\section*{Acknowledgements}
{\sloppy The research reported here was funded by the Leverhulme Trust, grant number RPG-2020-078, as part of the project “Declining case: inflectional\is{inflection} loss in progress” (\url{https://www.smg.surrey.ac.uk/projects/declining-case-inflectional-loss-in-progress}). The authors are grateful to two anonymous reviewers whose comments and suggestions led to significant improvements in the paper.\par}

%\section*{Contributions}
%John Doe contributed to conceptualization, methodology, and validation.
%Jane Doe contributed to writing of the original draft, review, and editing.


{\sloppy\printbibliography[heading=subbibliography,notkeyword=this]}
\end{document}
